% Generated from article.tex and its modular inputs by build.py.
% This file is independently compilable; edit the modular sources first.

% BEGIN preamble.tex
\documentclass[11pt]{article}
\usepackage[a4paper,margin=25mm]{geometry}
\usepackage[T1]{fontenc}
\usepackage{lmodern}
\usepackage{amsmath,amssymb,amsthm,mathtools}
\usepackage{booktabs,longtable,array,enumitem,xcolor,microtype,xurl}
\usepackage[colorlinks=true,linkcolor=blue!40!black,citecolor=blue!40!black,urlcolor=blue!50!black]{hyperref}
\usepackage{bookmark,fancyhdr}
\definecolor{researchblue}{RGB}{28,59,91}
\pagestyle{fancy}
\fancyhf{}
\fancyhead[L]{\small\color{researchblue}Polylogarithms, harmonic jets, and finite parts}
\fancyhead[R]{\small\thepage}
\renewcommand{\headrulewidth}{0.3pt}
\setlength{\headheight}{14pt}
\setlength{\emergencystretch}{2em}
\setcounter{tocdepth}{2}
\numberwithin{equation}{section}
\theoremstyle{plain}
\newtheorem{theorem}{Theorem}[section]
\newtheorem{lemma}[theorem]{Lemma}
\newtheorem{proposition}[theorem]{Proposition}
\newtheorem{corollary}[theorem]{Corollary}
\theoremstyle{definition}
\newtheorem{definition}[theorem]{Definition}
\newtheorem{example}[theorem]{Example}
\newtheorem{question}[theorem]{Research question}
\theoremstyle{remark}
\newtheorem{remark}[theorem]{Remark}
\DeclareMathOperator{\Li}{Li}
\DeclareMathOperator{\Log}{Log}
\DeclareMathOperator{\FP}{FP}
\DeclareMathOperator{\CT}{CT}
\DeclareMathOperator{\Res}{Res}
\DeclareMathOperator{\sgn}{sgn}
\newcommand{\C}{\mathbb C}
\newcommand{\R}{\mathbb R}
\newcommand{\N}{\mathbb N}
\newcommand{\Z}{\mathbb Z}
\newcommand{\Q}{\mathbb Q}
\newcommand{\ii}{\mathrm i}
\newcommand{\dd}{\,\mathrm d}
\newcommand{\src}[1]{\texttt{\small\nolinkurl{#1}}}
\newcommand{\zetad}[2]{\zeta^{(#1)}(#2)}
\newcommand{\fall}[2]{#1^{\underline{#2}}}
\setlist{topsep=5pt,itemsep=4pt,leftmargin=*}
\hypersetup{pdftitle={Mixed Resolvent Products, Higher Digamma Powers, and Stieltjes Collisions},pdfauthor={ProveIt research continuation}}

% END preamble.tex

\begin{document}
\hypersetup{pageanchor=false}
\begin{titlepage}
\thispagestyle{empty}
\vspace*{13mm}
{\color{researchblue}\Large ProveIt research continuation\par}
\vspace{15mm}
{\huge\bfseries Mixed Resolvent Products,\par
\vspace{3mm}Higher Digamma Powers,\par
\vspace{3mm}and Stieltjes Collisions\par}
\vspace{7mm}
{\Large Exact identities and normalized finite parts\par}
\vspace{12mm}
{\large 11 October 2026\quad$\cdot$\quad Version 1.0\par}
\vspace{10mm}
\begin{abstract}
\noindent
We develop three extensions of the current ProveIt polylogarithm and
Hurwitz--Stieltjes program. First, we obtain the Fourier spectrum and
Hurwitz transform of any finite product of cotangent resolvents with
independent complex powers. Opposite half-planes produce additional
endpoint poles; explicit full-amplitude subtraction gives a jointly
holomorphic germ whose coefficients are all prescribed Stieltjes,
argument, and exponent derivatives. An elementary Gamma specialization
exhibits genuinely path-dependent uncompleted limits. A finite
depth-two polylogarithmic formula evaluates the mixed logarithmic
transform at every Stieltjes index.

Second, we prove a terminating harmonic reduction and a normalized
antiderivative for every positive integer power of the digamma function.
We evaluate the fifth finite-part moment explicitly, retaining one
depth-four Laurent coefficient and a convergent depth-two derivative
combination, with all ordinary constants determined. The proof fixes
the entire Mittag--Leffler remainder and the finite asymptotic constants.

Third, we extend geometric collision subtraction to arbitrary Stieltjes
indices, ordinary argument derivatives, and analytic local weights.
Hermite interpolation removes the rational singular denominator before
the collision limit is taken. The limit is uniform with respect to
arbitrarily hierarchical gaps. The counterterms reduce to finite
generalized polylogarithms; three mixed cubic examples give explicit,
unequal order-zero anomalies.

We separate these proved statements from unresolved arithmetic
reductions, update the status of overlapping incoming questions, and
provide reproducible exact and numerical checks. No proof of the frozen
Gaussian $S_6$ or revised $S_8$ conjecture is claimed.
\end{abstract}
\vfill
\noindent\small
Source baseline: ProveIt revision
\src{9035bb0eaaf70f3398711b489733093dcc52e946}.\\
Prepared as an independently buildable research article with modular
\LaTeX{} sources, verification code, recorded results, and integration notes.
\end{titlepage}
\hypersetup{pageanchor=true}
\pagenumbering{roman}
\tableofcontents
\clearpage
\pagenumbering{arabic}

% BEGIN sections/01_scope_and_conventions.tex
\section{Scope, conventions, and the role of complete local subtraction}
\label{scope:section}

\subsection{Relation to the current manuscript and incoming reports}

The source baseline is the repository revision
\src{9035bb0eaaf70f3398711b489733093dcc52e946}, whose manuscript is
organized into four volumes and whose current incoming directory contains
seventeen ZIP reports \cite{repo:main}. The work here follows the
Hurwitz--Stieltjes and harmonic-identity spine. The source review included
the manuscript's structure and status records, the relevant finite-part
chapters, and the incoming reports bearing on the questions below.
It is a targeted mathematical audit, not a claim that every theorem in
the four volumes and all incoming archives has been re-proved.

There are already substantial results at this baseline. The cubic and
quartic digamma bridges, single-resolvent complex powers at nonpositive
resonances, all-degree formal Tornheim relation spaces, and the
Stieltjes-index-zero collision theorem are prior work
\cite{repo:cubic,repo:barnes,repo:quartic,repo:reflection}.
Several questions in older incoming reports have consequently become
stale. We preserve that distinction throughout: a theorem that was
already proved in a newer report is not presented as a new result here.

The concrete contributions and their limits are as follows.

\begin{longtable}{>{\raggedright\arraybackslash}p{0.24\textwidth}
                  >{\raggedright\arraybackslash}p{0.68\textwidth}}
\toprule
Target & Result in this article\\
\midrule
\endhead
Several resolvents with independent exponents &
Theorems~\ref{mr:fourier}, \ref{mr:hurwitz}, and \ref{mr:completion}
give the complete Fourier spectrum, the Hurwitz transform, and the
local completion for any finite number of nonreal parameters. Opposite
orientations are included. This answers the transformation and
normalization portions of Question 10 in \cite{repo:quartic}.\\[5pt]
Mixed logarithmic products &
Theorem~\ref{ml:allStieltjes} gives a finite expression of depth at
most two for two opposite logarithmic resolvents, at every Stieltjes
index and every ordinary argument-derivative order. It supplies a
specific instance of the mixed-product question in \cite{repo:quartic}.\\[5pt]
The fifth digamma moment &
Theorem~\ref{hp:fifth} resolves Question 1 of
\cite{repo:quartic}, with all finite corrections and a convergent
polygamma representation. It also supplies the first new case requested
by Question F3 of \cite{repo:barnes}.\\[5pt]
Arbitrary digamma powers &
Theorem~\ref{hp:allpowers} gives a terminating reduction to ordinary
convergent multiple zeta values and specified first Laurent coefficients.
Corollary~\ref{hp:primitive} fixes an antiderivative for every power.
No claim of an arithmetically minimal output basis is made.\\[5pt]
Higher-index geometric collisions &
Theorem~\ref{gcx:main} proves the total-collision limit for arbitrary
Stieltjes indices and ordinary derivatives, with an analytic local
weight and no restriction on gap ratios. This resolves the higher-index
extension posed in \cite{repo:reflection} at the level of continuity
after complete subtraction.\\[5pt]
Gaussian period conjectures &
The frozen $S_6$ and revised $S_8$ remain open. Neither a new period
certificate nor a counterexample is obtained. The earlier rejected
$S_8$ vector and the stated formal relation obstructions retain their
separate established status \cite{repo:main,repo:exact}.\\
\bottomrule
\end{longtable}

The word \emph{new} refers here to progress relative to this
pinned source corpus. Classical beta integrals, Gamma expansions,
harmonic products, and multiple-zeta dualities are cited as such.
The external literature check was focused; it does not establish global
priority for every consequence or specialization.

\subsection{Special-function and coefficient conventions}

The generalized Stieltjes functions are normalized by
\begin{equation}\label{scope:stieltjes}
 \zeta(1+u,x)=\frac1u+
       \sum_{m\ge0}\frac{(-1)^m}{m!}\gamma_m(x)u^m,
 \qquad \gamma_m=\gamma_m(1),\quad \gamma=\gamma_0.
\end{equation}
Thus $\gamma_0(x)=-\psi(x)$, where $\psi=\Gamma'/\Gamma$.
Equivalently,
\[
 \zeta(1-s,x)+\frac1s=\sum_{m\ge0}\gamma_m(x)\frac{s^m}{m!}.
\]
A superscript $(p)$ on $\gamma_m$ or $\psi$ means an ordinary derivative
in its argument. A prime on a zeta function of a single variable denotes
an order derivative unless an argument is explicitly specified.
The standard Hurwitz recurrence and differentiation formulas are
\cite{dlmf:hurwitz,dlmf:polygamma}
\begin{align}
 \gamma_m(x)&=\frac{\log^m x}{x}+\gamma_m(1+x),
       \label{scope:recurrence}\\
 \partial_x^p\zeta(s,x)&=(-1)^p(s)_p\zeta(s+p,x).
       \label{scope:hurwitzderivative}
\end{align}
Here $(s)_p=s(s+1)\cdots(s+p-1)$ is the rising factorial.
The symbol $[u^j]F$ means a Taylor or Laurent coefficient, including
no factorial. A derivative is recovered by multiplying a Taylor
coefficient by the appropriate factorial. This distinction matters at
every harmonic Laurent coefficient used below.

Our multiple zeta convention is strict and decreasing:
\[
 Z_d(s_1,\ldots,s_d)
   =\sum_{n_1>\cdots>n_d\ge1}
       n_1^{-s_1}\cdots n_d^{-s_d},
\]
in its convergence region, with meromorphic continuation where stated.
The notation $\zeta(s,t)$ in the polylogarithmic section denotes the
same depth-two function. For a nonnegative integer $N$,
$H_N^{(r)}=\sum_{n=1}^N n^{-r}$; multiple harmonic words use the same
strict ordering. No weakly ordered sum is silently substituted for a
strict one.

\subsection{Coordinate finite parts}

All endpoint finite parts are taken in the original integration
coordinate. For an integrand with a finite power--log singular expansion,
integrate after removing a small endpoint interval, expand the result
in the cutoff, and retain its constant term. At two distinct endpoints
the cutoffs are independent. In particular,
\begin{align}
 \FP\int_0^1 x^{\alpha-1}\log^j x\,\dd x
   &=\frac{(-1)^j j!}{\alpha^{j+1}}
                    &&(\alpha\ne0),\label{scope:monomial}\\
 \FP\int_0^1 x^{-1}\log^j x\,\dd x&=0
                    &&(j\ge0).\label{scope:resonantmonomial}
\end{align}
The first equality follows by meromorphic continuation from
$\Re\alpha>0$, or directly from the cutoff primitive.
At $\alpha=0$ the second equality is the coordinate constant;
it is not a value obtained by inserting zero into the first formula.

For periodic Stieltjes products, the local singularity occurs on the
right of each seam. We remove a right interval of ambient length
$\eta$ at every singular point and take independent constant terms.
Ordinary derivatives on the open intervals are used throughout.
Differentiating a periodic distribution that has already been extended
across a seam is a different operation and can carry additional
distributional contacts. Section~\ref{gcx:section} specifies its
convention explicitly.

For fixed distinct singular points, independent finite parts are linear
and commute with fixed parameter derivatives once enough local
monomials have been subtracted uniformly. This elementary fact is the
common analytic step behind the spectral and geometric constructions.
It does not authorize taking an arbitrary path through intersecting
parameter poles.

\begin{lemma}[Full-amplitude completion]\label{scope:completionlemma}
Let a parameter-dependent integrand near an endpoint have a finite
expansion whose singular terms are
$A_j(\mathbf t)x^{\alpha_j(\mathbf t)-1}$, with holomorphic amplitudes
and affine exponents, none identically zero. Suppose the subtracted remainder and each fixed
parameter derivative are uniformly integrable near $\mathbf t=0$.
For every $j$ such that $\alpha_j(0)=0$, subtract
$A_j(\mathbf t)/\alpha_j(\mathbf t)$ from its meromorphically continued
integral. The resulting local germ is holomorphic, and its Taylor
coefficients are exactly the coordinate finite parts of the
corresponding derivatives of the integrand. Power--log monomials follow
by differentiating in an auxiliary exponent.
\end{lemma}

\begin{proof}
On a fixed endpoint interval $(0,\rho)$, the integrated monomial is
\[
 A_j(\mathbf t)\frac{\rho^{\alpha_j(\mathbf t)}}
                            {\alpha_j(\mathbf t)}
 =\frac{A_j(\mathbf t)}{\alpha_j(\mathbf t)}
  +A_j(\mathbf t)
       \frac{\rho^{\alpha_j(\mathbf t)}-1}{\alpha_j(\mathbf t)}.
\]
The second term is entire in $\alpha_j$ and generates the cutoff
constants of all its logarithmic derivatives. Nonresonant monomials are
already holomorphic at the base point. Uniform integrability justifies
coefficient extraction from the remainder. The artificial splitting
point cancels against the rest of the integral.
\end{proof}

The amplitude in this lemma is the complete function
$A_j(\mathbf t)$, not merely its restriction to the polar hyperplane.
The latter specifies a residue but omits regular coefficients that
contribute to the prescribed finite part. The two-orientation example
in Section~\ref{mr:section} makes this distinction exact: different
uncompleted parameter paths give different finite limits for the same
specialized function.

% END sections/01_scope_and_conventions.tex


% BEGIN sections/02_resolvents.tex
\section{Products of resolvents with independent complex powers}
\label{mr:section}

Question 10 of the incoming \emph{Quartic Tails and Complex Resolvents}
report asks for products of several nonreal cotangent resolvents, including
parameters in opposite half-planes \cite{repo:quartic}. We answer its
analytic and transformation parts for any finite number of factors.
The answer has two distinct features: every Fourier coefficient has a
geometrically convergent Gamma-series representation, and the complete
endpoint subtraction depends on a signed sum of exponents as well as on
their total. The latter dependence produces genuine extra poles when both
half-planes occur.

\subsection{Branches and notation}

Let $z_1,\ldots,z_d\in\C\setminus\R$ and
$\boldsymbol\lambda=(\lambda_1,\ldots,\lambda_d)\in\C^d$. Repeated
values of the $z_j$ are allowed. Define
\begin{align}
 K(x)&=\pi\cot(\pi x),& R_z(x)&=(K(x)-z)^{-1},\notag\\
 \epsilon_j&=\sgn\Im z_j,&
 q_j&=\frac{z_j-\epsilon_j\pi\ii}{z_j+\epsilon_j\pi\ii},&
 C_j&=-\frac1{z_j+\epsilon_j\pi\ii},\label{mr:diskdata}\\
 a&=\sum_{\epsilon_j=1}\lambda_j,&
 b&=\sum_{\epsilon_j=-1}\lambda_j,&
 \Lambda&=a+b,\qquad E=a-b.\notag
\end{align}
In particular $|q_j|<1$. In every occurrence of $C_j^{\lambda_j}$ use
\[
 \Log C_j=\Log(1-q_j)-\log(2\pi)+\epsilon_j\pi\ii/2,
\]
where $\Log(1-q_j)$ is holomorphic on the disk and vanishes at $q_j=0$.
The continuous logarithm of $R_z$ on $(0,1)$
is fixed by $\Log R_z(x)-\log x\to0$ as $x\downarrow0$. For
$\epsilon=\sgn\Im z$, its imaginary part lies between $0$ and
$\epsilon\pi$. Powers in
\begin{equation}\label{mr:kernel}
 P_{\boldsymbol\lambda}(x)
       =\prod_{j=1}^d R_{z_j}(x)^{\lambda_j}
\end{equation}
use these logarithms. Put
\[
 B(\boldsymbol\lambda,\mathbf q)
       =\prod_j(1-q_j)^{\lambda_j},
 \qquad
 V(x;\boldsymbol\lambda)
       =\prod_j\bigl(\pi x\cot\pi x-z_jx\bigr)^{-\lambda_j}.
\]
The powers defining $B$ use the holomorphic logarithms on $|q_j|<1$.
The powers defining $V$ use their analytic logarithms at $x=0$, where
$V(0)=1$. Direct algebra gives both the global expression
\begin{align}
 P_{\boldsymbol\lambda}(x)
  ={}&\pi^{-\Lambda}B\,\sin^{\Lambda}(\pi x)e^{\pi\ii E x}
       \prod_{\epsilon_j=1}(1-q_je^{2\pi\ii x})^{-\lambda_j}
       \prod_{\epsilon_j=-1}(1-q_je^{-2\pi\ii x})^{-\lambda_j},
       \label{mr:sinefactor}
\end{align}
and the two endpoint expansions
\begin{equation}\label{mr:endpoints}
 P_{\boldsymbol\lambda}(x)=x^\Lambda V(x;\boldsymbol\lambda),\qquad
 P_{\boldsymbol\lambda}(1-y)
       =e^{\pi\ii E}y^\Lambda V(-y;\boldsymbol\lambda).
\end{equation}
The leading coefficient at zero is exactly one. This normalization
will fix every contact term below.

\subsection{An explicit Gamma series for the full Fourier spectrum}

Define the exponentially decreasing coefficient sequences
\begin{equation}\label{mr:UV}
 U_n=[w^n]\prod_{\epsilon_j=1}(1-q_jw)^{-\lambda_j},\qquad
 W_n=[w^n]\prod_{\epsilon_j=-1}(1-q_jw)^{-\lambda_j}.
\end{equation}
An empty product has coefficient $1$ at $n=0$ and $0$ otherwise.
Every coefficient is a finite sum; for example
\[
 U_n=\sum_{\substack{(n_j)\ge0\\\sum_{\epsilon_j=1}n_j=n}}
          \prod_{\epsilon_j=1}\frac{(\lambda_j)_{n_j}q_j^{n_j}}{n_j!}.
\]
For $k\in\Z$ set
\begin{equation}\label{mr:Gk}
 G_k(a,b)=
 \frac{e^{\pi\ii(a-b)/2}(-1)^k\Gamma(1+a+b)}
 {(2\pi)^{a+b}\Gamma(1+a-k)\Gamma(1+b+k)}.
\end{equation}
Reciprocal Gamma factors are read as entire functions. Restrictions
through an intersection of a numerator pole and a reciprocal-Gamma zero
are made as meromorphic restrictions, never as an informal $0\cdot\infty$.

\begin{theorem}[Fourier coefficients of any finite product]
\label{mr:fourier}
For $\Re\Lambda>-1$, the ordinary Fourier coefficients of
\eqref{mr:kernel} are
\begin{equation}\label{mr:ck}
 c_k(\boldsymbol\lambda,\mathbf z)
   :=\int_0^1P_{\boldsymbol\lambda}(x)e^{-2\pi\ii kx}\dd x
   =B\sum_{n,m\ge0}U_nW_mG_{k-n+m}(a,b).
\end{equation}
The double series converges normally, geometrically in the two summation
indices on compact subsets of $|q_j|<1$, away from its possible parameter
poles. It supplies a joint meromorphic continuation with possible poles
only on $\Lambda=-1,-2,\ldots$. Its mean is $c_0$.

If every $z_j$ is in the same half-plane, then
\begin{equation}\label{mr:samemean}
                         c_0=\prod_jC_j^{\lambda_j}.
\end{equation}
For mixed half-planes, \eqref{mr:ck} gives an explicit two-sided
spectrum and \eqref{mr:samemean} generally fails.
\end{theorem}

\begin{proof}
The classical Fourier beta integral \cite[Eq.~5.12.6]{dlmf:beta} reads
\[
 \int_0^1\sin^\Lambda(\pi x)e^{\pi\ii\nu x}\dd x
  =\frac{e^{\pi\ii\nu/2}\Gamma(1+\Lambda)}
         {2^\Lambda\Gamma(1+(\Lambda+\nu)/2)
                         \Gamma(1+(\Lambda-\nu)/2)},\qquad
 \Re\Lambda>-1.
\]
Substitution $\nu=a-b-2k$ gives $G_k$. Expanding only the analytic
factors in \eqref{mr:sinefactor}, whose series converge uniformly on
the circle, and integrating against the integrable endpoint factor
proves \eqref{mr:ck}.

Choose a radius $\rho>1$ with $\rho\max_j|q_j|<1$. Cauchy's
estimate bounds $U_n,W_n$ by compact-set constants times $\rho^{-n}$
and $\rho^{-m}$. Gamma reflection and the Gamma-ratio asymptotic show
that $G_k$ has at most polynomial growth in the integer $k$, uniformly
on compact parameter sets avoiding $\Lambda=-1,-2,\ldots$.
Consequently the series and all fixed parameter derivatives converge
normally there. Formula \eqref{mr:Gk} proves the stated meromorphy.

In a common half-plane put $w=e^{2\epsilon\pi\ii x}$. The kernel is
\[
 \left(\prod_jC_j^{\lambda_j}\right)
       (1-w)^\Lambda\prod_j(1-q_jw)^{-\lambda_j}.
\]
Its holomorphic disk function has constant coefficient $\prod C_j^{\lambda_j}$.
Abel limits are integrable when $\Re\Lambda>-1$, by the combined
endpoint exponent, and give \eqref{mr:samemean}. This argument uses
the total exponent, not separate inequalities on individual exponents.
\end{proof}

\begin{remark}[A finite hypergeometric specialization]
With one factor at $z_+=\pi\ii$ of exponent $a$ and one lower-half-plane
factor of exponent $b$ and disk coordinate $r$, \eqref{mr:ck} gives
\begin{equation}\label{mr:gaussmean}
 c_0=(1-r)^bG_0(a,b)\,
       {}_2F_1\left(\begin{matrix}-a,b\\1+b\end{matrix};r\right).
\end{equation}
This follows from $G_m/G_0=(-a)_m/(1+b)_m$ in the initial
nonexceptional domain. At exceptional parameters the full expression is
continued as a unit. Thus even one opposite orientation changes the
constant-mode formula from a product of constants to a hypergeometric
function. Formula \eqref{mr:ck}, rather than a name for a new function,
specifies the general answer by elementary finite coefficients and Gamma
values.
\end{remark}

\subsection{Hurwitz, Stieltjes, and log-Gamma transforms}

Let $p\ge0$ be an integer and put
\[
 A_p(s)=\prod_{j=1}^p(s-j),\quad A_0(s)=1,\qquad
 \sigma_+=-2\pi\ii,\quad \sigma_-=2\pi\ii,
 \quad \Log\sigma_\pm=\log(2\pi)\mp\pi\ii/2.
\]
Define, initially in an absolute-convergence region,
\begin{equation}\label{mr:Dirichlet}
 D_+(u)=\sum_{n\ge1}\frac{c_n}{n^u},\qquad
 D_-(u)=\sum_{n\ge1}\frac{c_{-n}}{n^u}.
\end{equation}
The definitions retain all coefficients, including their normalization $B$.

\begin{theorem}[The joint spectral identity]
\label{mr:hurwitz}
In the ordinary-integral domain
\[
 \Re\Lambda>-1,\qquad \Re(s+\Lambda)>p,
\]
away from the Hurwitz pole for $p=0$, analytic continuation from an
absolutely convergent Fourier region gives
\begin{align}
 J_p(s,\boldsymbol\lambda)
 &:=\int_0^1\partial_x^p\zeta(1-s,x)P_{\boldsymbol\lambda}(x)\dd x\notag\\
 &=\Gamma(s)\left\{\sigma_+^{p-s}D_+(s-p)
                      +\sigma_-^{p-s}D_-(s-p)\right\}.
       \label{mr:Hurwitz}
\end{align}
The completed Stieltjes generator is
\begin{equation}\label{mr:T}
 T_p(s,\boldsymbol\lambda)=J_p(s,\boldsymbol\lambda)
                                  +\mathbf1_{p=0}\frac{c_0}{s}.
\end{equation}
It is jointly meromorphic. Its only possible polar hypersurfaces are
\begin{equation}\label{mr:poles}
       s+\Lambda=p-j\quad(j\ge0),\qquad
       \Lambda=-1-j\quad(j\ge0).
\end{equation}
For $\Re\Lambda>p$,
\begin{equation}\label{mr:ordinaryStieltjes}
 \int_0^1\gamma_m^{(p)}(x)P_{\boldsymbol\lambda}(x)\dd x
                  =m![s^m]T_p(s,\boldsymbol\lambda).
\end{equation}
\end{theorem}

\begin{proof}
In an initial region $\Re s>p+1$, $\Re\Lambda>-1$, Fourier pairing
with the classical Hurwitz expansion \cite[Eq.~25.11.9]{dlmf:hurwitz}
gives zero for the constant mode and
\[
 \int_0^1\partial_x^p\zeta(1-s,x)e^{\pm2\pi\ii nx}\dd x
       =\Gamma(s)\sigma_\pm^{p-s}n^{p-s}.
\]
Here the differentiated Hurwitz series converges absolutely and
uniformly, and $P_{\boldsymbol\lambda}$ belongs to $L^1(0,1)$.
Thus termwise integration is justified directly by its uniform bound;
no regularity assertion for the endpoint values of $P$ is required.
This proves \eqref{mr:Hurwitz} in that region.

Use the exact holomorphic splitting
\begin{equation}\label{mr:split}
 \partial_x^p\left(\zeta(1-s,x)+\frac1s\right)
      =A_p(s)x^{s-p-1}+h_p(s,x),\qquad
 h_p(s,x)=\partial_x^p\left(\zeta(1-s,1+x)+\frac1s\right).
\end{equation}
It proves normal integrability on the stated ordinary domain and gives
analytic continuation of the Fourier identity to that domain. At the
two endpoints, the monomial exponents are respectively
$s+\Lambda-p-1+j$ and $\Lambda+j$. Subtracting any required finite
number of such monomials leaves a normally convergent integral.
Their integration denominators give exactly \eqref{mr:poles}.
The pole of the uncompleted Hurwitz factor is $-c_0/s$ for $p=0$;
\eqref{mr:T} removes it before Stieltjes coefficient extraction.
The convention
$\zeta(1-s,x)+s^{-1}=\sum_{m\ge0}\gamma_m(x)s^m/m!$
then gives \eqref{mr:ordinaryStieltjes}.
\end{proof}

The apparent coordinate poles of separate Dirichlet or Gamma factors
in \eqref{mr:Hurwitz} may cancel. Equation \eqref{mr:poles} concerns
the complete function, as proved by its endpoint integrals.

\begin{corollary}[An ordinary log-Gamma identity]
\label{mr:loggamma}
For $\Re\Lambda>-1$,
\begin{align}
 \int_0^1\log\Gamma(x)P_{\boldsymbol\lambda}(x)\dd x
 ={}&\frac{\log(2\pi)}2c_0\notag\\
 &+\sum_{\epsilon\in\{+,-\}}
     \frac{(\gamma+\Log\sigma_\epsilon)D_\epsilon(1)
                                     -D_\epsilon'(1)}{\sigma_\epsilon}.
       \label{mr:logGamma}
\end{align}
The prime differentiates the Dirichlet order. This is an ordinary
integral; its logarithmic endpoint is integrable under the stated
condition. The formula follows by differentiating \eqref{mr:Hurwitz}
at $s=1$ and using
$\zeta'(0,x)=\log\Gamma(x)-\tfrac12\log(2\pi)$.
\end{corollary}

For a concrete ordinary evaluation, take the special points
$z_\pm=\pm\pi\ii$ with exponents $2,1$. Their Fourier coefficients are
$\ii(-1,3,-3,1)/(8\pi^3)$ at modes $-1,0,1,2$, respectively.
Consequently \eqref{mr:logGamma} reduces to
\begin{align}
 &\int_0^1\log\Gamma(x)\,
           \frac{\sin^3(\pi x)}{\pi^3}e^{\pi\ii x}\dd x\notag\\
 &\qquad=
 \frac{3\{\gamma+\log(2\pi)\}-\log2}{32\pi^4}
       +\frac{\ii\{12\log(2\pi)-7\}}{64\pi^3}.
       \label{mr:loggammaexample}
\end{align}

\subsection{The full completion at each total-exponent resonance}

The preceding polar arrangement is completed explicitly as follows.
Set
\begin{equation}\label{mr:vandb}
 v_j(\boldsymbol\lambda)=[x^j]V(x;\boldsymbol\lambda),\qquad
 b_j(s,\boldsymbol\lambda)=[x^j]h_p(s,x)V(x;\boldsymbol\lambda).
\end{equation}
These are finite, effectively computable expressions. For instance
\[
 v_0=1,\qquad v_1=\sum_j\lambda_jz_j,\qquad
 v_2=\frac12\left(\sum_j\lambda_jz_j\right)^2
       +\frac12\sum_j\lambda_jz_j^2+\frac{\pi^2}3\Lambda.
\]
For $p+j>0$, the coefficient of $x^j$ in $h_p$ is
\[
 \frac{(-1)^{p+j}(1-s)_{p+j}}{j!}\zeta(1-s+p+j),
\]
and at $p=j=0$ it is $\zeta(1-s)+1/s$. Thus all required local
derivatives use only ordinary zeta derivatives and Stieltjes constants.

Fix $\boldsymbol\lambda^0$ with
$L=\sum_j\lambda_j^0\in\Z$, write
$\boldsymbol\lambda=\boldsymbol\lambda^0+\boldsymbol\eta$ and
$\delta=\sum_j\eta_j$. Define
\begin{align}
 H_{p,\boldsymbol\lambda^0}(s,\boldsymbol\eta)
 ={}&T_p(s,\boldsymbol\lambda^0+\boldsymbol\eta)
 -\mathbf1_{L\le p}\frac{A_p(s)v_{p-L}(\boldsymbol\lambda^0+\boldsymbol\eta)}
                                  {s+\delta}\notag\\
 &-\mathbf1_{L\le-1}
 \frac{\left\{1+(-1)^{-L-1}e^{\pi\ii E(\boldsymbol\lambda^0+\boldsymbol\eta)}\right\}
       b_{-L-1}(s,\boldsymbol\lambda^0+\boldsymbol\eta)}{\delta}.
       \label{mr:H}
\end{align}
False-indicator terms are omitted, not evaluated.

\begin{theorem}[All Stieltjes, argument, and exponent derivatives]
\label{mr:completion}
The function \eqref{mr:H} is holomorphic near $(s,\boldsymbol\eta)=(0,0)$.
For $m\ge0$ and a multi-index $\mathbf k\ge0$,
\begin{align}
 &\FP_x\int_0^1\gamma_m^{(p)}(x)
       \prod_j R_{z_j}(x)^{\lambda_j^0}
                       (\Log R_{z_j}(x))^{k_j}\dd x\notag\\
 &\hspace{20mm}=m!\left(\prod_j k_j!\right)
                    [s^m\boldsymbol\eta^{\mathbf k}]
                             H_{p,\boldsymbol\lambda^0}(s,\boldsymbol\eta).
       \label{mr:alljets}
\end{align}
The finite part uses the original endpoint coordinates $x$ and $1-x$.
The individual exponents $\lambda_j^0$ need not be integers; only
their total is integral in this statement. If the total is not integral,
the corresponding local germ of $T_p$ itself gives all the finite parts.
\end{theorem}

\begin{proof}
At zero, \eqref{mr:split} and \eqref{mr:endpoints} give the exact
families
\[
 A_p(s)v_j(\boldsymbol\lambda)x^{s+\Lambda-p-1+j},\qquad
 b_j(s,\boldsymbol\lambda)x^{\Lambda+j}.
\]
At the other endpoint the smooth Hurwitz factor is $h_p(s,-y)$.
Its product with the kernel has coefficient
$e^{\pi\ii E}(-1)^jb_j(s,\boldsymbol\lambda)$ at $y^{\Lambda+j}$.
The only exponents whose integrated denominators vanish at the base
point are $j=p-L$ and $j=-L-1$. Their complete amplitudes give the
two subtraction terms in \eqref{mr:H}.

For rigor, choose a fixed small endpoint interval $[0,\rho]$, and
subtract sufficiently many of the displayed monomials that the remainders
and their fixed derivatives are integrable uniformly on a parameter
polydisk. Every nonresonant integrated monomial is holomorphic there.
For each resonant exponent $\alpha-1$ use
\[
 \frac{\rho^\alpha}{\alpha}
             =\frac1\alpha+\frac{\rho^\alpha-1}{\alpha}.
\]
The second term is entire; its Taylor coefficients are exactly the
constant terms of the logarithmic cutoff integrals. Subtracting the
first term with its full parameter-dependent amplitude therefore makes
the remainder holomorphic and proves \eqref{mr:alljets}. Contributions
depending on the artificial splitting point cancel between the local
and remaining integrals. The same argument proves the nonintegral case,
where no denominator is resonant at the base point.
\end{proof}

\begin{corollary}[Why opposite half-planes produce an extra pole]
\label{mr:phasepole}
On the hyperplane $\Lambda=-M$, $M\ge1$, the regular-endpoint
residue in \eqref{mr:H} is
\begin{equation}\label{mr:mixedresidue}
       (1-e^{-2\pi\ii b})\,b_{M-1}(s,\boldsymbol\lambda).
\end{equation}
If all factors lie in the same half-plane this vanishes identically on
the hyperplane, so its quotient is removable. With both orientations
present it is generically nonzero. Its vanishing at an isolated integral
parameter vector does not make that multivariable quotient removable.
\end{corollary}
\begin{proof}
On $\Lambda=-M$, $E=\Lambda-2b=-M-2b$, and hence
$1+(-1)^{M-1}e^{\pi\ii E}=1-e^{-2\pi\ii b}$.
In a single orientation $b=0$ or $b=\Lambda$. With two orientations,
$b$ varies independently along the hyperplane. The local coefficient
$b_{M-1}$ may introduce further cancellations; generic nonvanishing
already occurs for $M=1$ and the constant test function.
\end{proof}

\subsection{An exactly computable mixed example}

For $z_+=\pi\ii$, $z_-=-\pi\ii$, write their exponents as
$\alpha,\beta$. The kernel and ordinary mean simplify to
\begin{align}
 P_{\alpha,\beta}(x)
   &=\left(\frac{\sin\pi x}{\pi}\right)^{\alpha+\beta}
                           e^{\pi\ii(\alpha-\beta)x},\label{mr:specialkernel}\\
 M(\alpha,\beta)
   &=\frac{e^{\pi\ii(\alpha-\beta)/2}\Gamma(1+\alpha+\beta)}
      {(2\pi)^{\alpha+\beta}\Gamma(1+\alpha)\Gamma(1+\beta)}.
       \label{mr:specialmean}
\end{align}
For any $c,d\in\C$ with $c+d\ne0$, continuation along
$\alpha=-1+ct$, $\beta=dt$ yields
\begin{equation}\label{mr:pathlimit}
           \lim_{t\to0}M(-1+ct,dt)=-2\pi\ii\frac{c}{c+d}.
\end{equation}
The fixed-coordinate finite part of the specialized integrand is instead
\begin{equation}\label{mr:fixedmean}
       \FP_x\int_0^1(K(x)-\pi\ii)\dd x=-\pi\ii.
\end{equation}
The exact correction is the full regular-endpoint quotient
\[
 \frac{1+e^{\pi\ii(\alpha-\beta)}}{1+\alpha+\beta}.
\]
Its limit on the same ray is $-\pi\ii(c-d)/(c+d)$; subtracting it
from \eqref{mr:pathlimit} gives \eqref{mr:fixedmean} for every ray.
Thus the rays $(c,d)=(1,0),(1,1),(0,1)$ yield continued limits
$-2\pi\ii,-\pi\ii,0$, respectively, for the same specialized
function. This is an exact counterexample to an uncompleted independent
exponent limit, not an error in the earlier single-resolvent theorem.

At a general negative-total hyperplane $\alpha+\beta=-M$, the
coefficient entering the mean residue is explicitly
\begin{equation}\label{mr:residuepoly}
 [x^{M-1}]V(x;\alpha,-M-\alpha)
          =(2\pi\ii)^{M-1}\binom{\alpha+M-1}{M-1}.
\end{equation}
Indeed this coefficient is the residue at zero of
$e^{\pi\ii(2\alpha+M)x}(\pi/\sin\pi x)^M\dd x$; putting
$w=e^{2\pi\ii x}$ converts it to
$(2\pi\ii)^{M-1}\Res_{w=1}w^{\alpha+M-1}(w-1)^{-M}\dd w$.
Consequently the residue of \eqref{mr:specialmean} is
\[
 (1-e^{2\pi\ii\alpha})(2\pi\ii)^{M-1}
                              \binom{\alpha+M-1}{M-1}.
\]
For example, at $\alpha=\beta=-1/2$ it is $2$, and direct
integration of $\pi\csc\pi x$ gives
\begin{equation}\label{mr:cscfp}
               \FP_x\int_0^1\pi\csc(\pi x)\dd x
                             =2\log(2/\pi).
\end{equation}

The dependence on the path can survive even when every transverse
linear path gives the expected zero. At $(-1,-1)$ write
\[
 \alpha=-1+(T+D)/2,\qquad \beta=-1+(T-D)/2.
\]
On the curved path $T=cD^2$ with $c\ne0$, the Gamma quotient
\eqref{mr:specialmean} has limit $\pi^2/c$. Indeed its leading terms
are $\pi^2D^2/T-\pi^2T$. The specialized kernel is
$\pi^2\csc^2\pi x$, whose coordinate finite part is zero.
Its full contact is
$(1-e^{\pi\ii D})\pi\ii D/T$, with the same curved-path limit.
The completed value is again independent of the path.

\begin{corollary}[A finite-part logarithmic family in ordinary constants]
\label{mr:negativeLog}
Let $R_\pm=R_{\pm\pi\ii}$ and put
\begin{align}
 Q(a,b)&=e^{-(a+b)\log(2\pi)+\pi\ii(a-b)/2}
              \frac{\Gamma(1+a+b)}{\Gamma(1+a)\Gamma(1+b)},\notag\\
 H(a,b)&=\frac{-2\pi\ii aQ(a,b)-1+e^{\pi\ii(a-b)}}{a+b}.
       \label{mr:negativeLoggerm}
\end{align}
This quotient is holomorphic at $(0,0)$. For integers $r,s\ge0$,
\begin{equation}\label{mr:negativeLogmoment}
 \FP_x\int_0^1(K-\pi\ii)(\Log R_+)^r(\Log R_-)^s\dd x
                      =r!s![a^rb^s]H(a,b).
\end{equation}
Every such value is a finite polynomial in $\pi\ii$, $\log(2\pi)$,
and ordinary zeta values. In particular,
\begin{align}
 \FP_x\int_0^1(K-\pi\ii)\Log R_+\dd x
       &=\frac{\pi^2}{2}+2\pi\ii\log(2\pi),\notag\\
 \FP_x\int_0^1(K-\pi\ii)(\Log R_-)^s\dd x
       &=\frac{(-\pi\ii)^{s+1}}{s+1}.
       \label{mr:negativeLogexamples}
\end{align}
\end{corollary}

\begin{proof}
At $\alpha=-1+a,\beta=b$, \eqref{mr:specialmean} equals
$-2\pi\ii aQ(a,b)/(a+b)$. Subtracting the complete paired endpoint
term gives \eqref{mr:negativeLoggerm}.
On $b=-a$, reflection gives
$Q(a,-a)=e^{\pi\ii a}\sin(\pi a)/(\pi a)$, so the numerator vanishes
identically. Analytic division and the endpoint proof above justify
all coefficients. The Gamma expansion proves the stated constant class.
For an independent check of the second identity,
$(\Log R_-)'=K-\pi\ii$. The endpoint logarithms are
$\log\eta+o(1)$ and $\log\eta-\pi\ii+o(1)$; the constant term of
their $(s+1)$st-power difference is $(-\pi\ii)^{s+1}$.
\end{proof}

\subsection{All polynomial moments and parameter derivatives}

The same beta formula gives an ordinary exponential moment at the
special points:
\begin{equation}\label{mr:expmean}
 \int_0^1e^{tx}P_{\alpha,\beta}(x)\dd x
 =\frac{e^{\pi\ii(\alpha-\beta)/2+t/2}\Gamma(1+\alpha+\beta)}
 {(2\pi)^{\alpha+\beta}
   \Gamma(1+\alpha+t/(2\pi\ii))
   \Gamma(1+\beta-t/(2\pi\ii))}.
\end{equation}
It is valid for $\Re(\alpha+\beta)>-1$ and every complex $t$.
Away from zeros of its value at $t=0$, its logarithmic derivatives are
\begin{align}
 \kappa_1&=\frac12+\frac{\psi(1+\beta)-\psi(1+\alpha)}{2\pi\ii},\notag\\
 \kappa_r&=-\frac{\psi^{(r-1)}(1+\alpha)
                    +(-1)^r\psi^{(r-1)}(1+\beta)}{(2\pi\ii)^r},
                    \qquad r\ge2.\label{mr:cumulants}
\end{align}
Thus its $n$th moment is $M(\alpha,\beta)Y_n(\kappa_1,\ldots,\kappa_n)$,
where $Y_n$ is the complete exponential Bell polynomial. At Gamma zeros,
differentiate the full reciprocal-Gamma formula \eqref{mr:expmean};
the factored Bell expression is only its meromorphic rewriting.
All mixed exponent derivatives follow by ordinary parameter
differentiation in this domain, and by \eqref{mr:H} at resonances.

At nonnegative integer $\alpha=m,\beta=n$, these formulas give finite
generalized-harmonic expressions at every moment order. Explicitly,
\[
 M(m,n)=
 \frac{e^{\pi\ii(m-n)/2}}{(2\pi)^{m+n}}\binom{m+n}{m},
 \qquad
 \kappa_1=\frac12+\frac{H_n-H_m}{2\pi\ii},
\]
and for $r\ge2$,
\[
 \kappa_r=
 \frac{(-1)^{r+1}(r-1)!}{(2\pi\ii)^r}
 \left\{\zeta(r)-H_m^{(r)}
                 +(-1)^r\bigl(\zeta(r)-H_n^{(r)}\bigr)\right\}.
\]
Substitution into the Bell polynomial is a finite identity involving
only harmonic numbers, zeta values, and the stated prefactor.

Finally, for any fixed-coordinate moment $\mathcal M$ in
\eqref{mr:alljets}, differentiation in a resolvent parameter gives
\[
 \partial_{z_j}\mathcal M(\boldsymbol\lambda,\mathbf k)
 =\lambda_j\mathcal M(\boldsymbol\lambda+\mathbf e_j,\mathbf k)
   +k_j\mathcal M(\boldsymbol\lambda+\mathbf e_j,\mathbf k-\mathbf e_j).
\]
The usual finite integration-by-parts formula in powers of $\Log R_{z_j}$
therefore supplies a normalized primitive in each $z_j$. For example,
with the other factors fixed and $\lambda_j\ne1$, a primitive at
$k_j=k$ is
\[
 \sum_{h=0}^k\frac{(-1)^hk!}{(k-h)!(\lambda_j-1)^{h+1}}
      \mathcal M(\boldsymbol\lambda-\mathbf e_j,
                         \mathbf k-h\mathbf e_j).
\]
When $\lambda_j=1$ it is the same moment with exponent $0$, logarithmic
degree $k+1$, divided by $k+1$. Subtract its value at a chosen base point
in the same half-plane to fix the additive constant. The proof is
coefficientwise differentiation of the completed integrals; no
uncompleted singular restriction is used.

% END sections/02_resolvents.tex


% BEGIN sections/03_polylogarithms.tex
\section{Mixed logarithmic moments in finite polylogarithmic depth}
\label{ml:section}

The general complex-power transform retains an explicit Dirichlet
coordinate. Its first mixed exponent derivative admits a substantially
smaller expression: two opposite logarithmic resolvents require only
depth-two polylogarithms, at every Stieltjes index and every argument
derivative. This section proves that reduction with its full finite-part
normalization.

\subsection{Two and three ordinary logarithmic moments}

Take $z$ in the upper half-plane and $w$ in the lower half-plane. Set
\[
 q=\frac{z-\pi\ii}{z+\pi\ii},\qquad
 r=\frac{w+\pi\ii}{w-\pi\ii},\qquad
 c_z=\Log\left(-\frac1{z+\pi\ii}\right),\quad
 c_w=\Log\left(-\frac1{w-\pi\ii}\right),
\]
with the branches fixed in Section~\ref{mr:section}. Define
\[
 L_z(x)=\Log R_z(x)-c_z,\qquad
 L_w(x)=\Log R_w(x)-c_w.
\]
Their $L^2$-convergent Fourier expansions are
\begin{equation}\label{ml:Fourierlog}
 L_z(x)=\sum_{n\ge1}\frac{q^n-1}{n}e^{2\pi\ii nx},\qquad
 L_w(x)=\sum_{n\ge1}\frac{r^n-1}{n}e^{-2\pi\ii nx}.
\end{equation}
In particular their means vanish.

\begin{proposition}[A dilogarithmic covariance]
\label{ml:dilog}
The ordinary integral is
\begin{equation}\label{ml:dilogeq}
 \int_0^1L_z(x)L_w(x)\dd x
       =\Li_2(qr)-\Li_2(q)-\Li_2(r)+\zeta(2)=:\mathcal D(q,r).
\end{equation}
Equivalently, the uncentered integral is $c_zc_w+\mathcal D(q,r)$.
For two parameters in the same half-plane, the centered integral is zero.
\end{proposition}
\begin{proof}
Pair the two $L^2$ Fourier series. Their constant product coefficient
is $\sum_{n\ge1}(q^n-1)(r^n-1)/n^2$, which is absolutely
convergent and sums to \eqref{ml:dilogeq}. Two equally oriented
series have no zero product mode. The logarithmic endpoint bounds
also verify ordinary integrability directly.
\end{proof}

We use the decreasing-index convention
\begin{equation}\label{ml:MPL}
 \Li_{s_1,\ldots,s_d}(u_1,\ldots,u_d)
 =\sum_{n_1>\cdots>n_d\ge1}
       \frac{u_1^{n_1}\cdots u_d^{n_d}}
            {n_1^{s_1}\cdots n_d^{s_d}}.
\end{equation}
Individual colors may have modulus greater than one; the actual coupled
summands, or equivalently their cumulative products, determine
convergence. Continuous limiting values are used when a displayed
quotient of colors is $0/0$.

For two upper-half-plane parameters with disk coordinates $q_1,q_2$
and one lower-half-plane parameter with coordinate $r$, let
$\nu_q(q)=1$, $\nu_q(1)=-1$. Then
\begin{align}
 &\int_0^1L_{z_1}(x)L_{z_2}(x)L_w(x)\dd x\notag\\
 &\quad=\sum_{a\in\{q_1,1\}}
        \sum_{b\in\{q_2,1\}}
        \sum_{c\in\{r,1\}}
        \nu_{q_1}(a)\nu_{q_2}(b)\nu_r(c)\notag\\
 &\hspace{24mm}\times
       \left\{\Li_{2,1}(ac,b/a)+\Li_{2,1}(bc,a/b)\right\}.
       \label{ml:triple}
\end{align}
This follows from the absolutely convergent elementary identity
\[
 \sum_{n,m\ge1}\frac{x^ny^m}{nm(n+m)}
     =\Li_{2,1}(x,y/x)+\Li_{2,1}(y,x/y),
\]
using $1/(nm(n+m))=1/((n+m)^2n)+1/((n+m)^2m)$.
The two cumulative colors are $ac,bc$, both of modulus at most one.
At three zero disk coordinates the value is $-2\zeta(3)$.
Formula \eqref{ml:triple} is thus a finite weight-three identity with
the branch and color conventions explicit.

\subsection{A depth-two transform with an arbitrary spectral order}

For $\Re u>0$, define
\begin{equation}\label{ml:Tu}
 \mathcal T_u(x,y)
       =\Li_{u+1}(x)\Li_1(xy)
                              -\Li_{1,u+1}(xy,1/y).
\end{equation}
For $x=y=1$ replace the right side, whose separate terms diverge, by
the explicitly convergent value
\begin{equation}\label{ml:T11}
       \mathcal T_u(1,1)=\zeta(u+1,1)+\zeta(u+2).
\end{equation}
Here $\zeta(s,t)=\sum_{n>m\ge1}n^{-s}m^{-t}$ is strict. At zero
colors the limit is used; in particular $\mathcal T_u(0,y)=
\mathcal T_u(x,0)=0$. Set
\begin{equation}\label{ml:Pu}
 \mathcal P_u(q,r)=\mathcal T_u(q,r)-\mathcal T_u(q,1)
                              -\mathcal T_u(1,r)+\mathcal T_u(1,1).
\end{equation}

\begin{lemma}[The two Fourier Dirichlet series]
\label{ml:coefficient}
If $f(x)=L_z(x)L_w(x)$, then
\begin{equation}\label{ml:Dirichlet}
 \sum_{k\ge1}\frac{\widehat f(k)}{k^u}=\mathcal P_u(q,r),\qquad
 \sum_{k\ge1}\frac{\widehat f(-k)}{k^u}=\mathcal P_u(r,q),
                          \qquad\Re u>0.
\end{equation}
All sums in this assertion converge absolutely.
\end{lemma}

\begin{proof}
The positive coefficient of the product is
\[
 \widehat f(k)=\sum_{m\ge1}
          \frac{(q^{m+k}-1)(r^m-1)}{m(m+k)},\qquad k\ge1.
\]
The four expanded sums are bounded by a constant multiple of
$\sum_{k\ge1}H_k/k^{\Re u+1}$, which converges for $\Re u>0$.
For $|x|,|y|\le1$ in the required cases,
\begin{align*}
 \sum_{k,m\ge1}\frac{x^{k+m}y^m}{k^um(k+m)}
 &=\sum_{k,m\ge1}\frac{x^k(xy)^m}{k^{u+1}m}
       -\sum_{n>k\ge1}\frac{(xy)^ny^{-k}}{nk^{u+1}}\\
 &=\mathcal T_u(x,y).
\end{align*}
The first equality uses
$1/(m(m+k))=k^{-1}(m^{-1}-(m+k)^{-1})$.
At $(x,y)=(1,1)$ the unseparated sum is
$\sum H_k/k^{u+1}=\zeta(u+1,1)+\zeta(u+2)$.
Expansion of the four numerator terms gives the first formula in
\eqref{ml:Dirichlet}; exchanging $q,r$ gives the second.
\end{proof}

The expressions in \eqref{ml:Tu} are finite combinations of depth at
most two. When $u$ is not integral, \eqref{ml:MPL} is an explicit
multiple Dirichlet sum with a complex order; spectral differentiation
inserts logarithms of its integer indices. Meromorphic continuation is
always taken for the complete combination, not for a divergent product
at $(x,y)=(1,1)$.

\subsection{All Stieltjes indices and all argument derivatives}

Write the local analytic logarithms
\[
 \ell_z(x)=-\Log(\pi x\cot\pi x-zx),\qquad
 \ell_w(x)=-\Log(\pi x\cot\pi x-wx),
\]
so that $L_z(x)=\log x-c_z+\ell_z(x)$ and similarly for $w$.
Define the finite coefficients $t_{p,h}$ by
\begin{equation}\label{ml:localt}
 [x^p]\{(X-c_z+\ell_z(x))(X-c_w+\ell_w(x))\}
                            =\sum_{h=0}^2t_{p,h}X^h.
\end{equation}
For example,
\begin{align*}
 (t_{0,2},t_{0,1},t_{0,0})&=(1,-c_z-c_w,c_zc_w),\\
 (t_{1,2},t_{1,1},t_{1,0})&=(0,z+w,-c_zw-c_wz).
\end{align*}
For $p\ge1$, $t_{p,2}=0$. Taylor expansion of the cotangent computes
every required coefficient exactly.

\begin{theorem}[Finite depth-two Stieltjes identity]
\label{ml:allStieltjes}
For $p\ge0$ the expression
\begin{align}
 \mathcal E_p(s;q,r)={}&\Gamma(s)
       \left\{\sigma_+^{p-s}\mathcal P_{s-p}(q,r)
                  +\sigma_-^{p-s}\mathcal P_{s-p}(r,q)\right\}
        +\mathbf1_{p=0}\frac{\mathcal D(q,r)}s\notag\\
 &\quad-A_p(s)\sum_{h=0}^2\frac{(-1)^hh!\,t_{p,h}}{s^{h+1}}
       \label{ml:E}
\end{align}
is holomorphic at $s=0$. For every $m\ge0$,
\begin{equation}\label{ml:allindex}
 \boxed{\displaystyle
 \FP_x\int_0^1\gamma_m^{(p)}(x)L_z(x)L_w(x)\dd x
                      =m![s^m]\mathcal E_p(s;q,r).}
\end{equation}
In particular the order-$m$ answer uses only finitely many order
derivatives of the depth-two functions in \eqref{ml:Tu}--\eqref{ml:Pu},
ordinary zeta functions, Gamma functions, logarithms, and the explicitly
finite local coefficients \eqref{ml:localt}.
\end{theorem}

\begin{proof}
For $\Re s>p+1$, pairing with the uniformly convergent Hurwitz Fourier
series and applying Lemma~\ref{ml:coefficient} gives the first line of
\eqref{ml:E}, before its last subtraction, as the integral of
$\partial_x^p(\zeta(1-s,x)+1/s)L_zL_w$. Its mean contribution is
exactly $\mathbf1_{p=0}\mathcal D/s$ by Proposition~\ref{ml:dilog}.

At zero, use \eqref{mr:split}. The smooth part times $L_zL_w$ is
locally integrable. Among the terms multiplying $A_p(s)x^{s-p-1}$,
only the coefficient $x^p$ has a resonant integration exponent.
By \eqref{ml:localt}, its integrated amplitude is
\[
 A_p(s)\sum_{h=0}^2t_{p,h}
                 \int_0^1x^{s-1}\log^hx\dd x
     =A_p(s)\sum_{h=0}^2\frac{(-1)^hh!t_{p,h}}{s^{h+1}}.
\]
At $x=1$ the Hurwitz factor is smooth and the logarithmic kernel is
integrable, so there is no second resonant term. The finite-subtraction
proof of Theorem~\ref{mr:completion} now applies verbatim to these
power-log monomials: removing their full amplitudes gives a holomorphic
germ whose coefficients are the fixed-coordinate finite parts.
This proves both assertions, including analytic continuation of the
complete depth-two expression wherever it is needed near $s=0$.
\end{proof}

For $p=0$ the subtraction is particularly transparent:
\begin{align}
 \mathcal E_0(s;q,r)={}&\Gamma(s)
   \left\{(-2\pi\ii)^{-s}\mathcal P_s(q,r)
                        +(2\pi\ii)^{-s}\mathcal P_s(r,q)\right\}
     +\frac{\mathcal D(q,r)}s\notag\\
 &\hspace{5mm}-\frac2{s^3}-\frac{c_z+c_w}{s^2}-\frac{c_zc_w}s.
       \label{ml:pzero}
\end{align}
Keeping only the leading cubic pole would lose two prescribed endpoint
constants. For $p>0$, keeping only pole residues would additionally lose
the regular coefficients contributed by the polynomial $A_p(s)$.

\subsection{A special-point identity involving a harmonic Laurent coefficient}

Let $L=\log(2\pi)$ and define $\eta$ by the strict expansion
\begin{equation}\label{ml:eta}
 \zeta(1+s,1)=\frac1{s^2}+\frac\gamma s
       +\frac{\gamma^2-\zeta(2)}2+\eta s+O(s^2).
\end{equation}
This is the same coefficient used in the cubic and quartic incoming
reports \cite{repo:cubic,repo:barnes}. At $q=r=0$,
$\mathcal P_s(0,0)=\zeta(1+s,1)+\zeta(2+s)$ and
$\mathcal D(0,0)=\zeta(2)$. The kernel is real:
\[
 L_{\pi\ii}(x)L_{-\pi\ii}(x)
      =\log^2(2\sin\pi x)+\pi^2(x-\tfrac12)^2.
\]
Expanding \eqref{ml:pzero} gives the exact identity
\begin{align}
 &\FP_x\int_0^1\gamma_0(x)
       \left\{\log^2(2\sin\pi x)+\pi^2(x-\tfrac12)^2\right\}\dd x
       \notag\\
 &\quad=2\eta+2\zeta'(2)-\frac{\gamma^3+L^3}{3}
                -\gamma\zeta(2)-\frac L2\zeta(2)-\frac23\zeta(3).
       \label{ml:specialStieltjes}
\end{align}
For completeness, put $A=\gamma+L$. The first four coefficients of
$\Gamma(s)(2\pi)^{-s}\cos(\pi s/2)$ are
\[
 \frac1s-A+\left(\frac{A^2}{2}-\frac{\zeta(2)}4\right)s
    +\left(-\frac{A^3}{6}+\frac{A\zeta(2)}4-\frac{\zeta(3)}3\right)s^2.
\]
Multiplication by \eqref{ml:eta} plus $\zeta(2+s)$ proves
\eqref{ml:specialStieltjes}. This calculation also provides a direct
connection with the harmonic-jet power identities developed below.
It does not claim an ordinary-constant evaluation of $\eta$.

\subsection{A generating identity for every mixed logarithmic moment}

At the special resolvent points, the Gamma mean \eqref{mr:specialmean}
gives a compact all-order ordinary identity. For
$F_+=L_{\pi\ii}$ and $F_-=L_{-\pi\ii}$,
\begin{equation}\label{ml:jointcumulants}
 \int_0^1e^{aF_+(x)+bF_-(x)}\dd x
       =\frac{\Gamma(1+a+b)}{\Gamma(1+a)\Gamma(1+b)}.
\end{equation}
Initially $\Re(a+b)>-1$; in particular it is analytic at $(0,0)$.
Taking its logarithm gives
\begin{equation}\label{ml:logjoint}
 \sum_{n\ge2}\frac{(-1)^n\zeta(n)}n
                     \left((a+b)^n-a^n-b^n\right).
\end{equation}
This power series is used locally, for example when
$|a|,|b|,|a+b|<1$.
Consequently every mixed cumulant with $p,q\ge1$ is
\begin{equation}\label{ml:cumulant}
              (-1)^{p+q}(p+q-1)!\zeta(p+q),
\end{equation}
and the pure centered cumulants vanish. Exponentiating
\eqref{ml:logjoint} is a finite coefficient algorithm for every moment.
For example,
\begin{align}
 \int F_+F_-&=\zeta(2),&
 \int F_+^2F_-&=-2\zeta(3),\notag\\
 \int F_+^2F_-^2&=11\zeta(4),&
 \int F_+^3F_-^2&=-24\zeta(5)-12\zeta(2)\zeta(3).
       \label{ml:momentexamples}
\end{align}
All integrals in \eqref{ml:momentexamples} are ordinary integrals over
$(0,1)$. This is a beta/Gamma specialization, with the classical Gamma
expansion retaining its attribution \cite{dlmf:gamma}; its role here is
to evaluate the two-orientation logarithmic moment family exactly.

\subsection{Derivatives and primitives in two resolvent parameters}

Since $\partial_zc_z=C_z$, differentiation gives
$\partial_zL_z=R_z-C_z$. Thus \eqref{ml:allindex} also supplies a
two-parameter primitive for moments of centered resolvents. For the
uncentered version put
\[
 S_{m,p}(z)=\FP_x\int_0^1\gamma_m^{(p)}(x)\Log R_z(x)\dd x,
\]
whose single-logarithm formula is already proved in
\cite[complex-power section]{repo:quartic}. Because
$\FP\int_0^1\gamma_m^{(p)}=0$, the uncentered mixed-logarithm moment is
\begin{equation}\label{ml:uncentered}
 \mathcal U_{m,p}(z,w)=m![s^m]\mathcal E_p(s;q,r)
                                      +c_zS_{m,p}(w)+c_wS_{m,p}(z).
\end{equation}
It satisfies the exact differential identity
\begin{equation}\label{ml:doubleprimitive}
 \partial_z\partial_w\mathcal U_{m,p}(z,w)
       =\FP_x\int_0^1\gamma_m^{(p)}(x)R_z(x)R_w(x)\dd x.
\end{equation}
Integrating over a rectangle in the product of the two half-planes
therefore yields the alternating four evaluations of
$\mathcal U_{m,p}$ at its corners, which fixes all integration constants.
Partial fractions give the additional consistency check
$R_zR_w=(R_z-R_w)/(z-w)$. The differential identity is valid for
the completed coordinate finite parts and does not presume that separate
uncompleted spectral limits can be multiplied.

% END sections/03_polylogarithms.tex


% BEGIN sections/04_higher_digamma.tex
\section{A finite harmonic reduction for every digamma power}
\label{hp:sec:all}

The newest source reports contain two independent fourth-power bridges.
The report \emph{Quartic Tails and Complex Resolvents}, Question~1 in its
research section, asks for the fifth power with every finite correction\cite{repo:quartic};
Question~2 asks about polynomials that remove lower-depth data.
\emph{Barnes Resolvents and Higher Harmonic Identities}, Question~F3, asks for
all higher powers and explicitly allows additional harmonic numerators\cite{repo:barnes}.
We answer the fifth-power problem and give a terminating reduction for
every fixed power. The latter has a precisely specified output: ordinary
convergent multiple zeta values and first Laurent coefficients of strict
multiple zeta functions. It does not assert that those coefficients
reduce to ordinary Stieltjes constants.

Write
\[
 H_N(s_1,\ldots,s_d)
   =\sum_{N\ge n_1>\cdots>n_d\ge1}
             n_1^{-s_1}\cdots n_d^{-s_d},\qquad
 Z_d(s_1,\ldots,s_d)=\lim_{N\to\infty}H_N(s_1,\ldots,s_d)
\]
when the latter converges. The larger index is first. The empty word
has value one, and $H_N^{(r)}=H_N(r)$, $H_N=H_N^{(1)}$.
The classical harmonic products, Euler-sum reductions and digamma
expansions keep their existing attribution\cite{dlmf:psi,hoffman,mezo}.
All indices in a harmonic word below are positive integers, except
for the one explicitly moving outer order. We use
\[
 \zeta(1+u)=u^{-1}+\sum_{j\ge0}\frac{(-1)^j\gamma_j}{j!}u^j,
 \qquad \gamma=\gamma_0,
 \qquad Q_k=\FP\int_0^1\psi(x)^k\,\dd x.
\]
Here $\FP$ is the constant in the cutoff coordinate $x=\varepsilon$:
it is the coefficient of $\varepsilon^0(\log\varepsilon)^0$ in
$\int_\varepsilon^1$. The upper endpoint is regular. Consequently
\[
 \FP\int_0^1x^{-j}\,\dd x=-\frac1{j-1}\quad(j\ge2),
 \qquad \FP\int_0^1x^{-1}\,\dd x=0.
\]

\subsection{The coefficients and the uniform theorem}

For $n\ge0$ introduce the formal series
\begin{equation}\label{hp:V}
 V_n(t)=-1+(H_n-\gamma)t
           +\sum_{r\ge2}\bigl(H_n^{(r)}+(-1)^r\zeta(r)\bigr)t^r.
\end{equation}
Only finitely many terms are needed in every operation that follows.
For $1\le j\le k$ put
\begin{equation}\label{hp:pdef}
 p_{k,j}(n)=[t^{k-j}]V_n(t)^k,
 \qquad \kappa_k=[t^k]V_0(t)^k,
 \qquad \Delta p_{k,j}(n)=p_{k,j}(n)-p_{k,j}(n-1).
\end{equation}
Define
\begin{align}
 B_k={}&\kappa_k+
   \sum_{j=2}^k\left\{
    \frac1{j-1}\sum_{n\ge1}\frac{\Delta p_{k,j}(n)}{n^{j-1}}
        -\sum_{n\ge1}\frac{p_{k,j}(n)}{n^j}\right\},
       \label{hp:Bdef}\\
 C_k={}&[L^0]\sum_{n=1}^N\frac{p_{k,1}(n)}n,
       \qquad L=\log N,\label{hp:Cdef}\\
 E_k(u)={}&\sum_{n\ge1}\Delta p_{k,1}(n)n^{-u}
       \qquad(\Re u>0).\label{hp:Edef}
\end{align}
The notation in \eqref{hp:Cdef} means the constant in the polynomial
asymptotic expansion in $L$, with an $o(1)$ remainder; existence is
part of the theorem. If $k=1$, the sum in \eqref{hp:Bdef} is empty.

\begin{theorem}[Uniform finite harmonic reduction]\label{hp:allpowers}
For every integer $k\ge1$, all series in \eqref{hp:Bdef} converge
absolutely, $C_k$ exists, $E_k$ is meromorphic near zero, and
\begin{equation}\label{hp:universal}
 \boxed{\displaystyle Q_k=B_k-C_k+[u^1]E_k(u).}
\end{equation}
Every term on the right has a finite algebraic construction from
\eqref{hp:V}. More precisely:
\begin{enumerate}
 \item $B_k$ and $C_k$ are finite linear combinations of convergent
 multiple zeta values, with coefficients in
 $\Q[\gamma,\zeta(2),\ldots,\zeta(k)]$. No limiting harmonic sum
 remains after the finite reduction given below.
 \item $E_k(u)$ is a finite linear combination of functions
 $Z_{d+1}(m+u,s_1,\ldots,s_d)$, where all displayed fixed indices are
 positive integers, $m+s_1+\cdots+s_d\le k-1$, and $d+1\le k-1$.
 Coefficients are in the same ordinary-constant algebra.
 \item For $k\ge2$, its only term of depth $k-1$ is
 \begin{equation}\label{hp:topdepth}
   -k!\,Z_{k-1}(1+u,\underbrace{1,\ldots,1}_{k-2}).
 \end{equation}
 All other moving-order terms have depth at most $k-2$.
\end{enumerate}
For $k=1$, $E_1=0$ and $Q_1=0$. For $k=2$, the formula gives
$Q_2=2\gamma_1-2\zeta(2)$.
\end{theorem}

The depth statement describes the specified finite normal form.
It is not an arithmetic independence statement. In particular,
shuffle, stuffle, duality, or an additional analytic identity may
further reduce the retained coordinates.

\subsection{The entire remainder and the normalized primitive}

\begin{lemma}[Mittag--Leffler identity at every power]\label{hp:ML}
Let $P_{k,n}(z)=\sum_{j=1}^kp_{k,j}(n)(n+z)^{-j}$. Away from the
nonpositive integers,
\begin{equation}\label{hp:MLformula}
 \psi(z)^k=P_{k,0}(z)+\kappa_k+
                    \sum_{n\ge1}\bigl(P_{k,n}(z)-P_{k,n}(0)\bigr).
\end{equation}
The series converges normally on compact subsets of that domain.
\end{lemma}

\begin{proof}
The digamma recurrence and its Taylor expansion at zero give
$t\psi(-n+t)=V_n(t)$. Therefore \eqref{hp:pdef} gives exactly the
principal parts, and $\kappa_k$ is the constant at the origin.
To determine the entire remainder, set $f(z)=\psi(z)^k-P_{k,0}(z)$
and integrate $zf(w)/(w(w-z))$ on $|w|=N+1/2$.
Uniformly on those circles, $f(w)=O(\log^kN)$.
In the right half-plane this is the usual digamma asymptotic.
In the left half-plane use
$\psi(w)=\psi(1-w)-\pi\cot(\pi w)$. The cotangent is uniformly
bounded on the chosen circles: where $|\Im w|\ge1/4$ its exponential
formula gives the bound; on the remaining arcs $\Re w$ stays a fixed
distance from the integers for sufficiently large $N$.
The contour integral is consequently $O(\log^kN/N)$ and tends to zero.
Its residues give \eqref{hp:MLformula}, including the constant
$f(0)=\kappa_k$.

Finally $p_{k,j}(n)=O(1+\log^{k-j}n)$ and
$(n+z)^{-j}-n^{-j}=O(n^{-j-1})$ on a fixed compact set.
This proves normal convergence. Matching principal parts alone
would not have proved the lemma: the vanishing contour integral
is what fixes the entire function.
\end{proof}

For $n\ge1$ set
\[
 F_1(n,z)=\Log(1+z/n)-z/n,\qquad
 F_j(n,z)=\frac{(n+z)^{1-j}-n^{1-j}}{1-j}-\frac{z}{n^j}
 \quad(j\ge2).
\]
On $\Re z>0$ use the logarithm real on the positive real axis.
The lemma immediately gives the following all-power antiderivative.

\begin{corollary}[A primitive with fixed additive constant]\label{hp:primitive}
For every $k\ge1$,
\begin{equation}\label{hp:primitiveformula}
 \mathcal P_k(z)=p_{k,1}(0)\Log z+
      \sum_{j=2}^k\frac{p_{k,j}(0)z^{1-j}}{1-j}
      +\kappa_kz+
      \sum_{n\ge1}\sum_{j=1}^kp_{k,j}(n)F_j(n,z)
\end{equation}
satisfies $\mathcal P_k'(z)=\psi(z)^k$ on $\Re z>0$.
Its expansion at the singular origin has constant zero. Thus
$\mathcal P_k(b)=\FP\int_0^b\psi(x)^k\,\dd x$ for $b>0$, and
ordinary integrals between positive endpoints are differences of
$\mathcal P_k$. Every fixed argument derivative of
\eqref{hp:MLformula} is also justified by normal convergence.
\end{corollary}

\begin{proof}
The displayed $F_j$ are the primitives, vanishing at zero, of
$(n+z)^{-j}-n^{-j}$. Their size on a fixed compact set is
$O(n^{-j-1})$; the same normal-convergence argument applies to
the differentiated series. On $|z|<1$ all nonsingular summands
vanish at zero. This proves the normalization as well as the
derivative identity.
\end{proof}

\subsection{Proof of the finite reduction}

Integrate \eqref{hp:MLformula}. For each $j\ge2$ finite summation by
parts and the vanishing terminal boundary term give
\begin{align*}
 &-\frac{p_{k,j}(0)}{j-1}
   +\sum_{n\ge1}p_{k,j}(n)
       \left\{\frac{n^{1-j}-(n+1)^{1-j}}{j-1}-n^{-j}\right\}\\
 &\hspace{15mm}=
   \frac1{j-1}\sum_{n\ge1}\frac{\Delta p_{k,j}(n)}{n^{j-1}}
         -\sum_{n\ge1}\frac{p_{k,j}(n)}{n^j}.
\end{align*}
All these sums are absolutely convergent. Hence
\begin{equation}\label{hp:afterintegration}
 Q_k=B_k+\sum_{n\ge1}p_{k,1}(n)
             \left\{\log(1+1/n)-1/n\right\}.
\end{equation}
Write $p(n)=p_{k,1}(n)$ and $d(n)=p(n)-p(n-1)$. At a finite cutoff,
\begin{equation}\label{hp:finiteparts}
 \sum_{n=1}^Np(n)\left\{\log(1+1/n)-1/n\right\}
 =p(N)\log(N+1)-\sum_{n=1}^Nd(n)\log n
                         -\sum_{n=1}^N\frac{p(n)}n.
\end{equation}
No divergent infinite rearrangement has been used.

The harmonic asymptotics imply, for a polynomial $P$,
$p(N)=P(\log N)+O(N^{-1}\log^M N)$. Thus the first term on the
right of \eqref{hp:finiteparts} has zero constant in $L=\log N$.
Likewise
\begin{equation}\label{hp:Dasym}
 d(n)=\frac{\sum_{r=0}^{k-2}a_r\log^r n}{n}
            +O(n^{-2}\log^M n)
\end{equation}
for some coefficients $a_r$ (the case of a zero polynomial is allowed).
Consequently
\[
 E_k(u)=\sum_ra_r(-1)^r\zeta^{(r)}(1+u)
   +\sum_{n\ge1}\left(d(n)-\frac{\sum_ra_r\log^r n}{n}\right)n^{-u},
\]
where the second sum is holomorphic for $\Re u>-1$, locally normally
with every fixed derivative. The Stieltjes coefficient convention gives
\begin{equation}\label{hp:coefficientconstant}
 [u^1]E_k(u)=-[L^0]\sum_{n=1}^Nd(n)\log n.
\end{equation}
Indeed $[u^1](-1)^r\zeta^{(r)}(1+u)=-\gamma_{r+1}$, while the
remaining coefficient is minus its absolutely convergent logarithmic
sum. Since \eqref{hp:afterintegration} converges, taking constants in
\eqref{hp:finiteparts} proves \eqref{hp:universal}.

Here are finite rules that compute every object in that proof. Expand
each harmonic polynomial by the strict quasi-shuffle rule
\begin{equation}\label{hp:stuffle}
 H_N(a)H_N(b,\mathbf v)
   =\sum_{\mathbf w\in(a)*(b,\mathbf v)}H_N(\mathbf w),
\end{equation}
where the recursive product inserts either initial letter or merges
them to their sum; repeated words retain their multiplicities. The
merge has a plus sign. For any nonempty word $\mathbf w=(w_1,\ldots,w_d)$,
\begin{align}
 \sum_{n=1}^N\frac{H_{n-1}(\mathbf w)}{n^j}
     &=H_N(j,\mathbf w),\label{hp:previous}\\
 \sum_{n=1}^N\frac{H_n(\mathbf w)}{n^j}
     &=H_N(j,\mathbf w)+H_N(j+w_1,w_2,\ldots,w_d).
       \label{hp:inclusive}
\end{align}
For the empty word the second term in \eqref{hp:inclusive} is absent.
The differences $\Delta p_{k,j}$ are computed by replacing
$H_n^{(r)}$ by $H_{n-1}^{(r)}+n^{-r}$ and subtracting.
Equations \eqref{hp:previous}--\eqref{hp:inclusive} then reduce $B_k$
to convergent multiple zeta values, because its initial index is
always at least two. They reduce $E_k$ to the finite moving-order
combination in the theorem.

For completeness, leading ones in $C_k$ can also be removed by an
explicit terminating rule. Let $\mathcal R_{\mathbf v}(T)$ denote
the harmonic regularization polynomial, with
\[
 \mathcal R_\varnothing(T)=1,\quad \mathcal R_{(1)}(T)=T,
 \quad\mathcal R_{\mathbf v}(T)=Z(\mathbf v)\quad(v_1\ge2).
\]
For $\mathbf v=(1,\mathbf w)$ with $a$ initial ones, write
$(1)*\mathbf w=a\mathbf v+\sum_{\mathbf u\ne\mathbf v}
m_{\mathbf u}\mathbf u$. Then
\begin{equation}\label{hp:regularization}
 \mathcal R_{\mathbf v}(T)=\frac1a
  \left\{T\mathcal R_{\mathbf w}(T)
       -\sum_{\mathbf u\ne\mathbf v}m_{\mathbf u}
                         \mathcal R_{\mathbf u}(T)\right\}.
\end{equation}
This decreases either total weight or the number of initial ones,
so it terminates. Before passing to $\mathcal R$, the same recurrence
with each terminal $Z(\mathbf v)$ replaced by $H_N(\mathbf v)$
is an exact finite identity, with $T=H_N$. Only afterward are the
convergent terminal words replaced by their limits. Their tails are
$O(N^{-1}\log^M N)$, and $H_N=L+\gamma+O(N^{-1})$.
Thus their constant in $L$ is obtained by $T=\gamma$.
This proves existence and the finite construction of $C_k$.
It also proves the individual continuation of every moving-order term:
when $m=1$, its inner harmonic word is a polynomial in $\log n$
plus $O(n^{-1}\log^M n)$, so a finite sum of zeta derivatives
subtracts its pole. When $m\ge2$, its defining series is already
holomorphic near $u=0$.

The index-weight bound follows directly from \eqref{hp:V}. Finally,
the unique harmonic monomial of degree $k-1$ in the simple-pole
coefficient is $-kH_n^{k-1}$. Its difference contributes
$-k(k-1)H_{n-1}^{k-2}/n$. The coefficient of the strict word of
$k-2$ ones in $H_{n-1}^{k-2}$ is $(k-2)!$.
Every other term has smaller harmonic degree after taking a
difference. This proves \eqref{hp:topdepth} and completes the theorem.

\section{The fifth-power identity and its additional spectral coordinate}
\label{hp:sec:fifth}

Define the inherited strict coefficients and the new height-one
coefficient by
\begin{equation}\label{hp:heightcoeffs}
 \eta=[u^1]Z_2(1+u,1),\qquad
 \theta=[u^1]Z_3(1+u,1,1),\qquad
 \phi=[u^1]Z_4(1+u,1,1,1).
\end{equation}
Brackets denote Laurent coefficients, with neither an additional
factorial nor a sign. A convergent depth-two derivative coordinate is
\begin{align}
 \chi_4&=2(\partial_1-\partial_2)Z_2(3,1)
                         +\partial_1Z_2(2,2)\label{hp:chi}\\
 &=-\sum_{n>m\ge1}\left\{
       \frac{2\log(n/m)}{n^3m}+\frac{\log n}{n^2m^2}\right\}.
       \notag
\end{align}
Both sums converge absolutely. The symbols $\phi$ and $\chi_4$ specify
the exact data retained below. There is no assertion that either is
arithmetically independent of the other constants.

\begin{theorem}[Fifth harmonic bridge]\label{hp:fifth}
The coordinate finite part of the fifth digamma power is
\begin{align}\label{hp:fifthmain}
 \boxed{\begin{aligned}
 Q_5={}&30\gamma^2\zeta(2)-60\gamma\zeta(3)+\frac{35}{3}\zeta(4)
       -120\phi+120\gamma\theta
       +60\bigl(\zeta(2)-\gamma^2\bigr)\eta\\
 &-20\bigl(\gamma^3-3\gamma\zeta(2)+2\zeta(3)\bigr)\gamma_1
       +10\chi_4-20\zeta'(4)
       +10\zeta(2)\zeta'(2)-20\gamma\zeta'(3).
 \end{aligned}}
\end{align}
An ordinary absolutely convergent series for the same constant is
\begin{equation}\label{hp:fifthseries}
 \boxed{\begin{aligned}
 Q_5={}&30\gamma^2\zeta(2)-60\gamma\zeta(3)+\frac{35}{3}\zeta(4)\\
      &+20\gamma_4-120\zeta(2)\gamma_2
       -5\sum_{n\ge1}\frac{\mathcal U(n)\log n}{n}.
 \end{aligned}}
\end{equation}
where, with $p=\psi(n)$ and $\ell=\log n$,
\begin{equation}\label{hp:Un}
 \mathcal U(n)=-4p^3+24\zeta(2)p-12p\psi'(n)-2\psi''(n)
    +4\ell^3-24\zeta(2)\ell
    +\frac{4\zeta(2)-2\psi'(n)}n+\frac{4p}{n^2}.
\end{equation}
The summand in \eqref{hp:fifthseries} is
$O(n^{-2}\log^3n)$.
\end{theorem}

\subsection{Explicit principal parts and every finite constant}

Put
\[
 b_n=H_n-\gamma,\quad c_n=H_n^{(2)}+\zeta(2),\quad
 d_n=H_n^{(3)}-\zeta(3),\quad e_n=H_n^{(4)}+\zeta(4).
\]
The principal coefficients supplied by \eqref{hp:pdef} are
\begin{equation}\label{hp:fifthparts}
 \begin{aligned}
 p_{5,5}&=-1,\qquad p_{5,4}=5b_n,\qquad p_{5,3}=5c_n-10b_n^2,\\
 p_{5,2}&=10b_n^3-20b_nc_n+5d_n,\qquad p_{5,1}=5r_n,\\
 r_n&=-b_n^4+6b_n^2c_n-2c_n^2-4b_nd_n+e_n.
 \end{aligned}
\end{equation}
The constant at zero is
\begin{equation}\label{hp:kappa5}
 \kappa_5=-\gamma^5+20\gamma^3\zeta(2)-30\gamma^2\zeta(3)
 -30\gamma\zeta(2)^2+20\gamma\zeta(4)
 +20\zeta(2)\zeta(3)-5\zeta(5).
\end{equation}
The exact difference of the residue is particularly simple:
\begin{equation}\label{hp:resdiff}
 r_n-r_{n-1}=\frac{4a_{n-1}}n+\frac{2c_{n-1}}{n^2}
                         +\frac{4b_{n-1}}{n^3},
 \qquad a_n=-b_n^3+3b_nc_n-d_n.
\end{equation}
It follows by substituting $b\mapsto b+n^{-1}$,
$c\mapsto c+n^{-2}$, $d\mapsto d+n^{-3}$, and
$e\mapsto e+n^{-4}$ into $r$. In particular, it is an exact finite
identity before any analytic continuation.

The lower-pole constant in \eqref{hp:Bdef} is
\begin{align}\label{hp:B5}
 B_5={}&-\gamma^5+30\gamma^3\zeta(2)+30\gamma^2\zeta(2)
      -80\gamma^2\zeta(3)-60\gamma\zeta(3)
      -\frac{65}{2}\gamma\zeta(4)\notag\\
 &\hspace{17mm}+60\zeta(2)\zeta(3)+\frac{35}{3}\zeta(4)-69\zeta(5).
\end{align}
We provide finite bookkeeping for this reduction and for the second
constant, so that the result is not dependent on a numerical fit.
The convergent values of weights at most five needed after the
quasi-shuffle expansion are
\begin{align}\label{hp:MZVtable}
 Z_2(2,1)&=\zeta(3),&Z_2(3,1)&=\tfrac14\zeta(4),&
 Z_2(2,2)&=\tfrac34\zeta(4),\notag\\
 Z_3(2,1,1)&=\zeta(4),&Z_4(2,1,1,1)&=\zeta(5),\notag\\
 Z_2(4,1)=Z_3(3,1,1)&=2\zeta(5)-\zeta(2)\zeta(3),\notag\\
 Z_2(3,2)=Z_3(2,2,1)&=3\zeta(2)\zeta(3)-\tfrac{11}{2}\zeta(5),\notag\\
 Z_2(2,3)=Z_3(2,1,2)&=\tfrac92\zeta(5)-2\zeta(2)\zeta(3).
\end{align}
These are classical Euler-sum and duality identities. For example,
Euler's formula for $Z_2(n,1)$ gives the $(4,1)$ value; the stuffle
product for $\zeta(2)\zeta(3)$ and its iterated-integral shuffle
product
\[
 \zeta(2)\zeta(3)=Z_2(2,3)+3Z_2(3,2)+6Z_2(4,1)
\]
give the two remaining depth-two values. Reversing the iterated
integral word gives the three indicated dualities. The height-one
values follow from the ordinary logarithmic beta integral. Also
$\zeta(2)^2=5\zeta(4)/2$.

For an additional check on \eqref{hp:B5}, let $B_{5,j}$ denote its
summand indexed by $j$ in \eqref{hp:Bdef}, excluding $\kappa_5$.
The finite rules above give
\begin{align}\label{hp:Bpieces}
 B_{5,2}={}&10\gamma^3\zeta(2)+30\gamma^2\zeta(2)
 -60\gamma^2\zeta(3)-70\gamma\zeta(3)+\tfrac{85}{2}\gamma\zeta(4)
 \notag\\
 &\quad+65\zeta(2)\zeta(3)+15\zeta(4)-\tfrac{215}{2}\zeta(5),\notag\\
 B_{5,3}={}&10\gamma^2\zeta(3)+10\gamma\zeta(3)-25\gamma\zeta(4)
       -30\zeta(2)\zeta(3)-5\zeta(4)+\tfrac{115}{2}\zeta(5),\notag\\
 B_{5,4}={}&5\gamma\zeta(4)+5\zeta(2)\zeta(3)
                            +\tfrac53\zeta(4)-15\zeta(5),\notag\\
 B_{5,5}={}&\zeta(5).
\end{align}
Adding these to \eqref{hp:kappa5} proves \eqref{hp:B5}.

To specify the constant of the simple-pole sum, set
\[
 A=4\gamma^3-12\gamma\zeta(2)+4\zeta(3),\quad
 C_0=-\gamma^4+6\gamma^2\zeta(2)-4\gamma\zeta(3)
                       -2\zeta(2)^2+\zeta(4),
 \quad e_j(N)=H_N(\underbrace{1,\ldots,1}_{j}).
\]
Expanding $r_n$ in strict words and applying
\eqref{hp:inclusive} gives the exact finite identity
\begin{align}\label{hp:finiteR}
 \sum_{n=1}^N\frac{r_n}{n}={}&-24e_5(N)+24\gamma e_4(N)
       +12(\zeta(2)-\gamma^2)e_3(N)+Ae_2(N)+C_0H_N\notag\\
 &-24H_N(2,1,1,1)+24\gamma H_N(2,1,1)
       +12(\zeta(2)-\gamma^2)H_N(2,1)+AH_N^{(2)}\notag\\
 &+4H_N(1,1,3)+2H_N(1,2,2)+4H_N(1,3,1)\notag\\
 &+4H_N(2,3)+2H_N(3,2)+4H_N(4,1)\notag\\
 &-4\gamma\{H_N(1,3)+H_N^{(4)}\}
          +2\zeta(2)\{H_N(1,2)+H_N^{(3)}\}.
\end{align}
One can apply \eqref{hp:regularization} directly. To make its
asymptotic evaluation especially short, write $T=H_N$ and note
\begin{align*}
 H_N(1,3)&=T\zeta(3)-\tfrac54\zeta(4)+o(1),\\
 H_N(1,2)&=T\zeta(2)-2\zeta(3)+o(1),\\
 4H_N(1,1,3)+2H_N(1,2,2)+4H_N(1,3,1)
   &=2T^2\zeta(3)-\tfrac52T\zeta(4)
                                  -2\zeta(2)\zeta(3)+o(1).
\end{align*}
Every $o(1)$ remains so after the displayed polynomial logarithmic
factors. The last line follows by stuffling $H_N$ into
$H_N(1,3)$, $H_N(2,2)$ and $H_N(3,1)$ and then using
\eqref{hp:MZVtable}. The elementary symmetric polynomials satisfy
\[
 e_j(N)=[t^j]\exp\left(Tt+
           \sum_{r=2}^j\frac{(-1)^{r-1}H_N^{(r)}}r t^r\right).
\]
Substitute these formulas and $T=L+\gamma+o(1)$ into
\eqref{hp:finiteR}. The complete result is
\begin{equation}\label{hp:Casym}
 \sum_{n=1}^N\frac{r_n}{n}
   =-\frac{L^5}{5}+4\zeta(2)L^3-18\zeta(4)L+C_r+o(1),
\end{equation}
where
\begin{equation}\label{hp:Cr}
 C_r=-\frac{\gamma^5}{5}+6\gamma^3\zeta(2)
      -16\gamma^2\zeta(3)-\frac{13}{2}\gamma\zeta(4)
      +12\zeta(2)\zeta(3)-\frac{69}{5}\zeta(5).
\end{equation}
Since $p_{5,1}=5r$, $C_5=5C_r$. Thus the substantial ordinary-constant
cancellation is exactly
\begin{equation}\label{hp:BCcancel}
 B_5-C_5=30\gamma^2\zeta(2)-60\gamma\zeta(3)+\frac{35}{3}\zeta(4).
\end{equation}
In particular, neither the finite constant at infinity nor the
finite parts of the higher poles have been discarded.

\subsection{The moving-order identity and the proof of the bridge}

Let $E_r(u)=E_5(u)/5$. From \eqref{hp:resdiff}, finite strict
harmonic identities give the full meromorphic germ
\begin{align}\label{hp:Er}
 E_r(u)={}&-24Z_4(1+u,1,1,1)+24\gamma Z_3(1+u,1,1)
         +12(\zeta(2)-\gamma^2)Z_2(1+u,1)\notag\\
 &+4(\gamma^3-3\gamma\zeta(2)+\zeta(3))\zeta(1+u)
       +4Z_2(1+u,3)\notag\\
 &+2Z_2(2+u,2)+2\zeta(2)\zeta(2+u)
                    +4Z_2(3+u,1)-4\gamma\zeta(3+u).
\end{align}
For example the identity used for its first line is
\[
 a_n=-6H_n(1,1,1)+6\gamma H_n(1,1)
       +3(\zeta(2)-\gamma^2)H_n
       +\gamma^3-3\gamma\zeta(2)+\zeta(3)+H_n^{(3)}.
\]
The series proof of \eqref{hp:Er} first applies in $\Re u>0$ and
then continues to the germ at zero. It is therefore legitimate to
extract its linear coefficient.

The only apparently singular additional depth-two term is removed
by an identity valid with a free outer order:
\[
 Z_2(1+u,3)=\zeta(1+u)\zeta(3)-\zeta(4+u)-Z_2(3,1+u).
\]
Consequently
\[
 [u^1]Z_2(1+u,3)
       =-\zeta(3)\gamma_1-\zeta'(4)-\partial_2Z_2(3,1).
\]
All remaining additional terms in \eqref{hp:Er} are ordinary
convergent derivatives. Extracting the coefficient and using
\eqref{hp:BCcancel} in \eqref{hp:universal} proves
\eqref{hp:fifthmain}.

For the independently useful convergent form, the elementary
polygamma relations give
\[
 b_{n-1}=\psi(n),\qquad c_{n-1}=2\zeta(2)-\psi'(n),\qquad
 d_{n-1}=\tfrac12\psi''(n).
\]
Substitution into \eqref{hp:resdiff} shows that
\[
 r_n-r_{n-1}
   =\frac{-4\log^3n+24\zeta(2)\log n+\mathcal U(n)}n.
\]
The standard positive-axis asymptotics of $\psi,\psi',\psi''$
give $\mathcal U(n)=O(n^{-1}\log^2 n)$.
The Stieltjes finite limits are
\[
 [L^0]\sum_{n=1}^N\frac{\log^j n}{n}=\gamma_j.
\]
Using them in \eqref{hp:coefficientconstant} proves
\eqref{hp:fifthseries}, with ordinary absolute convergence as
claimed.

\subsection{The highest-depth coefficient and a finite cancellation}

The normalization of $\phi$ has a direct elementary Mellin
description. In $\Re s>1$,
\begin{equation}\label{hp:heightMellin}
 Z_4(s,1,1,1)=\frac1{6\Gamma(s)}
   \int_0^\infty\frac{t^{s-1}[-\log(1-e^{-t})]^3}{e^t-1}\,\dd t.
\end{equation}
Subtract $(-\log t)^3/(6t)$ on $(0,1)$. The remainder is
holomorphic at $s=1$, and its value there is zero, because the
unsubtracted kernel is $-1/24$ times the derivative of
$[-\log(1-e^{-t})]^4$. Let $g_5=[u^5]\Gamma(1+u)^{-1}$.
Then
\begin{align}\label{hp:phiIntegral}
 \phi={}&g_5+\frac16\int_0^1\log t
  \left\{\frac{[-\log(1-e^{-t})]^3}{e^t-1}
                      -\frac{(-\log t)^3}{t}\right\}\,\dd t\notag\\
 &+\frac16\int_1^\infty
             \frac{\log t\,[-\log(1-e^{-t})]^3}{e^t-1}\,\dd t,
\end{align}
where both integrals converge ordinarily, and
\[
 g_5=\frac{\gamma^5-10\gamma^3\zeta(2)+20\gamma^2\zeta(3)
        +15\gamma\zeta(2)^2-30\gamma\zeta(4)
        -20\zeta(2)\zeta(3)+24\zeta(5)}{120}.
\]
This fixes the depth-four coordinate independently of a formal
Laurent label. It does not evaluate it in a smaller constant algebra.

The previously proved cubic and quartic identities\cite{repo:cubic,repo:quartic} are
\begin{align*}
 Q_3&=3\zeta(2)-6\eta-6\gamma\gamma_1,\\
 Q_4&=12\zeta(3)-12\gamma\zeta(2)+18\zeta(4)+4\zeta'(3)
       -24\theta+24\gamma\eta
       +12(\gamma^2-\zeta(2))\gamma_1.
\end{align*}
Combining them with the new theorem eliminates both inherited
strict coefficients:
\begin{corollary}[Eliminating the inherited height-one data]\label{hp:cancellation}
For $P_5(T)=T^5+5\gamma T^4+10(\gamma^2+\zeta(2))T^3$,
\begin{align}\label{hp:polyCancellation}
 \FP\int_0^1P_5(\psi(x))\,\dd x
 ={}&-120\phi+10\chi_4+
              \left(\frac{260}{3}+90\gamma\right)\zeta(4)
           -20\zeta'(4)+10\zeta(2)\zeta'(2)\notag\\
 &-20\bigl(\gamma^3+3\gamma\zeta(2)+2\zeta(3)\bigr)\gamma_1.
\end{align}
\end{corollary}
The extra derivative coordinate $\chi_4$ remains. Thus this is a
proved cancellation of $\eta$ and $\theta$, not a claimed solution
of the stronger all-power, single-coordinate polynomial problem.

\subsection{Verification and the exact scope of the result}

The accompanying code records sixty exact finite checks. It expands \eqref{hp:V} exactly, verifies
\eqref{hp:resdiff} and \eqref{hp:finiteR} on finite harmonic sums,
replays the terminating quasi-shuffle regularization, and reproduces
\eqref{hp:B5}--\eqref{hp:BCcancel}. Its numerical comparison uses
two analytically different methods: local integration of the Taylor
series of $(x\psi(x))^5$ followed by ordinary quadrature, and
\eqref{hp:fifthseries} with a finite asymptotic tail summed through
Hurwitz-zeta order derivatives. The direct evaluator gives
\[
 Q_5=-8.69030396262453826629558966236565357677029354492388429266\ldots.
\]
The recorded run agrees through forty-eight decimal places.
It is a floating-point diagnostic, not an interval certificate;
the exact proof is \eqref{hp:universal} together with the finite
coefficient calculations above. No numerical relation search is
used to infer a constant reduction.

The general theorem settles the existence of a finite reduction
for every power in its stated multiple-zeta-jet class and fixes
all entire and endpoint constants. The fifth-power theorem
supplies the explicit case requested in both newest reports.
The arithmetic reduction of $\chi_4$ or $\phi$, a proof that a
polynomial in $\psi$ can remove all lower-depth data at arbitrary
order, and a common-Hurwitz-shift counterpart remain separate
research problems.

% END sections/04_higher_digamma.tex


% BEGIN sections/05_higher_collisions.tex
\section{Polylogarithmic counterterms for higher Stieltjes collisions}
\label{gcx:section}

The geometric collision theorem at Stieltjes index zero in the incoming
\emph{Higher Reflection Identities and Resonant Zeta Calculus}\cite{repo:reflection} leaves two explicit gaps: logarithmic
Stieltjes singularities, and weighted collision formulas.  The theorem below
covers both, together with arbitrary ordinary argument derivatives.  The
counterterm is finite and explicitly reducible to iterated polylogarithms.
The proof also explains why no comparison between different shrinking gaps
is needed.  It does not evaluate the global merged cubic constant, and it
does not assert associativity of a different, successively chosen
subtraction prescription.

\subsection{Local singular polynomials and the finite-part convention}

Normalize generalized Stieltjes functions by
\[
 \zeta(1+u,a)=\frac1u+
       \sum_{m\ge0}\frac{(-1)^m}{m!}\gamma_m(a)u^m.
\]
For $m,r\ge0$, put
\begin{align}
 f_{m,r}(x)&=\gamma_m^{(r)}(\{x\}),\qquad
 h_{m,r}(z)=\gamma_m^{(r)}(1+z),\label{gcx:fg}\\
 d_r&=r+1,\qquad
 L_{m,r}(X)=\prod_{j=1}^{r}(\partial_X-j)X^m.
 \label{gcx:sing-polynomial}
\end{align}
The empty product is the identity.  Thus $L_{m,0}(X)=X^m$ and
$L_{0,r}(X)=(-1)^r r!$.  The exact one-sided local decomposition is
\begin{equation}\label{gcx:local}
 f_{m,r}(x)=\boldsymbol 1_{x>0}\,
              x^{-r-1}L_{m,r}(\log x)+h_{m,r}(x),
 \qquad -1<x<1,\quad x\ne0.
\end{equation}
Derivatives in this formula are ordinary derivatives on the open
intervals.  They are not derivatives of already extended periodic
distributions.

For completeness, the singular polynomial has the harmonic form
\begin{equation}\label{gcx:harmonic-poly}
 L_{m,r}(X)=(-1)^r r!
 \sum_{j=0}^{\min(m,r)}(-1)^j\binom mj j!
 e_j\left(1,\frac12,\ldots,\frac1r\right)X^{m-j},
\end{equation}
where $e_j$ is an elementary symmetric polynomial and $e_0=1$.
Indeed,
\[
 L_{m,r}(X)=(-1)^r m![z^m](1-z)_r e^{zX}.
\]
All analytic tails in \eqref{gcx:local} are holomorphic near zero.  Their
Taylor coefficients require only ordinary Stieltjes constants and zeta
jets: for $k\ge1$,
\begin{equation}\label{gcx:tail-derivative}
 \gamma_m^{(k)}(1)=(-1)^{m+k}m![u^m](1+u)_k\zeta(k+1+u).
\end{equation}
This follows by differentiating the Hurwitz identity
$\partial_a^k\zeta(s,a)=(-1)^k(s)_k\zeta(s+k,a)$.

Fix a periodic smooth weight $w$ which is real analytic in a neighborhood
of $0$; $w=1$ and $w(x)=e^{2\pi i n x}$ are allowed.  Fix indices
$(m_i,r_i)$, $1\le i\le d$.  For distinct sufficiently small real shifts
$s_i$, define
\begin{equation}\label{gcx:CQ}
 C_{\boldsymbol m,\boldsymbol r}^{w}(\boldsymbol s)
 =\FP\int_{\R/\Z} w(x)\prod_{i=1}^d f_{m_i,r_i}(x+s_i)\,\dd x,
 \qquad
 Q_{\boldsymbol m,\boldsymbol r}^{w}
 =\FP\int_0^1w(x)\prod_{i=1}^d\gamma_{m_i}^{(r_i)}(x)\,\dd x.
\end{equation}
At each singular point $a$, the finite part removes $(a,a+\eta)$ and
retains the constant term in the ambient coordinate $\eta$.  Distinct
points use independent cutoffs.  The same convention is used at the
merged endpoint.  In particular, no cutoff is divided by a shrinking
separation.

\subsection{A finite Hermite construction}

For a nonempty subset $S\subseteq\{1,\ldots,d\}$, write
\begin{align}
 d_i&=r_i+1,\quad N_S=\sum_{i\in S}d_i,\quad
 b_S=\min_{i\in S}s_i,\notag\\
 q_S(x;\boldsymbol s)&=w(x)\prod_{j\notin S}h_{m_j,r_j}(x+s_j),
 \qquad
 D_S(x;\boldsymbol s)=\prod_{i\in S}(x+s_i)^{d_i},
 \label{gcx:qD}\\
 \Lambda_S(x;\boldsymbol s)&=
              \prod_{i\in S}L_{m_i,r_i}(\log(x+s_i)),
 \qquad x>-b_S.
 \label{gcx:Lambda}
\end{align}
The polynomial $D_S(x;\boldsymbol s)$ is monic of degree $N_S$.  Let $P_S(x;\boldsymbol s)$ be the unique polynomial of degree
less than $N_S$ that agrees with $q_S$ at $x=-s_i$ through derivative
order $d_i-1$, for every $i\in S$.

\begin{definition}[Canonical logarithmic collision counterterm]
\label{gcx:counterterm-definition}
For $|S|\ge2$, define
\begin{equation}\label{gcx:T}
 T_S(\boldsymbol s)=\FP_{R\to\infty}\FP_{\eta\downarrow0}
 \int_{-b_S+\eta}^{R}
       \frac{P_S(x;\boldsymbol s)}{D_S(x;\boldsymbol s)}
                  \Lambda_S(x;\boldsymbol s)\,\dd x,
 \qquad
 T(\boldsymbol s)=\sum_{|S|\ge2}T_S(\boldsymbol s).
\end{equation}
The lower finite part is taken first, with the ambient displacement
$\eta$.  At infinity, discard the nonconstant polynomial in $\log R$
and retain its constant term.  These are two separate endpoint
operations, not a diagonal spectral regularization.
\end{definition}

The upper operation in \eqref{gcx:T} is explicit.  Put
\[
 p_S(\boldsymbol s)=[x^{N_S-1}]P_S(x;\boldsymbol s),\qquad
 \lambda_S(X)=\prod_{i\in S}L_{m_i,r_i}(X),\qquad
 \Psi_S(X)=\int_0^X\lambda_S(t)\,\dd t.
\]
The only divergent upper term is
$p_S(\boldsymbol s)\Psi_S(\log R)$.  The difference between its derivative
and the integrand is $O(R^{-2}(1+\log R)^{\sum_{i\in S}m_i})$.
For a singleton $S$, the same definition gives zero, since translation
to its sole endpoint leaves a sum of pure monomials
$x^{-k}(\log x)^j$, $k\ge1$, whose two-endpoint finite parts vanish.

\begin{theorem}[All indices, all ordinary derivatives, and analytic weights]
\label{gcx:main}
For the convention \eqref{gcx:CQ}, the counterterm \eqref{gcx:T} satisfies
\begin{equation}\label{gcx:limit}
 C_{\boldsymbol m,\boldsymbol r}^{w}(\boldsymbol s)
       -T(\boldsymbol s)
       \longrightarrow Q_{\boldsymbol m,\boldsymbol r}^{w}
 \qquad (\boldsymbol s\longrightarrow0)
\end{equation}
through distinct real shifts.  The convergence imposes no lower bound on
ratios of the gaps.  In particular, it holds for arbitrary hierarchical
collisions and for paths whose relative gaps vanish faster than every
power of their overall size.

Every $T_S$ is a finite combination of rational functions of the gaps,
values and derivatives of $w$ and the analytic Stieltjes tails at the
shifted nodes, and generalized polylogarithms of weight at most
$1+\sum_{i\in S}m_i$.  All logarithms in its defining real integral have
positive arguments.  If every $m_i=0$, the formula is exactly the
previous rational--logarithmic collision counterterm, including its
ordinary-derivative extension.
\end{theorem}

\begin{proof}[Proof of the collision limit]
Choose $\delta>0$ so small that $w$ and all required tails are holomorphic
on $|z|<3\delta$, and restrict $|s_i|<\delta/8$.  Use the circle coordinate
$(-\delta,1-\delta)$ and separate the local interval
$(-\delta,\delta)$.  The outside integral is continuous in the shifts
(and analytic when the outside weight is analytic).  The local product
expands by the subset $S$ of factors contributing their singular terms in
\eqref{gcx:local}.  The empty subset is analytic.  The nonempty subset is
\begin{equation}\label{gcx:local-subset}
 \FP\int_{-b_S}^{\delta}
        \frac{q_S(x;\boldsymbol s)}{D_S(x;\boldsymbol s)}
                  \Lambda_S(x;\boldsymbol s)\,\dd x.
\end{equation}
Only the endpoint $-b_S$ is singular in this integral: all other nodes
$-s_i$ of the subset lie to its left.

Both $P_S$ and the divided remainder
\[
 A_S(x;\boldsymbol s)=\frac{q_S(x;\boldsymbol s)-P_S(x;\boldsymbol s)}
                           {D_S(x;\boldsymbol s)}
\]
extend holomorphically through coalescing nodes.  This assertion is not
based on individually singular Lagrange coefficients.  It follows from
the fixed-contour formulas
\begin{align}
 P_S(x;\boldsymbol s)&=\frac1{2\pi i}\oint_{|z|=2\delta}
 \frac{q_S(z;\boldsymbol s)}{D_S(z;\boldsymbol s)}
 \frac{D_S(z;\boldsymbol s)-D_S(x;\boldsymbol s)}{z-x}\,\dd z,
 \label{gcx:Hermite-contour}\\
 A_S(x;\boldsymbol s)&=\frac1{2\pi i}\oint_{|z|=2\delta}
 \frac{q_S(z;\boldsymbol s)}{D_S(z;\boldsymbol s)(z-x)}\,\dd z,
 \qquad |x|<\delta.
 \label{gcx:remainder-contour}
\end{align}
These also give a uniform bound for $A_S$ on the local integration
interval after decreasing $\delta$ or the shift neighborhood slightly.

The finite-part tail of the principal part has the convergent form
\begin{align}
 B_S(\boldsymbol s):={}&\FP\int_\delta^\infty
      \frac{P_S(x;\boldsymbol s)}{D_S(x;\boldsymbol s)}
                  \Lambda_S(x;\boldsymbol s)\,\dd x\notag\\
 ={}&-p_S(\boldsymbol s)\Psi_S(\log\delta)
 +\int_\delta^\infty\left\{
      \frac{P_S(x;\boldsymbol s)}{D_S(x;\boldsymbol s)}
                    \Lambda_S(x;\boldsymbol s)
       -\frac{p_S(\boldsymbol s)}x\lambda_S(\log x)\right\}\,\dd x.
 \label{gcx:tail}
\end{align}
For $x\ge\delta$ use the holomorphic branches of $\Log(x+s_i)$ which are
real at real shifts.  The integrand and its fixed parameter derivatives
are locally uniformly $O(x^{-2}(1+\log x)^M)$, with
$M=\sum_{i\in S}m_i$.  Consequently $B_S$ is holomorphic near the full
diagonal.

Subtracting $T_S$ from \eqref{gcx:local-subset} leaves exactly
\begin{equation}\label{gcx:after-subtraction}
 \int_{-b_S}^{\delta}A_S(x;\boldsymbol s)
                \Lambda_S(x;\boldsymbol s)\,\dd x-B_S(\boldsymbol s).
\end{equation}
The rational denominators have disappeared from the remaining local
integral.  To prove the required uniformity, set $y=x+b_S$ and
$a_i=s_i-b_S\ge0$.  Its integrand becomes
\[
 A_S(y-b_S;\boldsymbol s)
       \prod_{i\in S}L_{m_i,r_i}(\log(y+a_i)).
\]
For $\delta$ small, all $y+a_i$ in this integral are less than $1$.
Because $y+a_i\ge y$, its absolute value is bounded by
$K(1+|\log y|)^M$, independently of every gap ratio.  This is integrable
at $y=0$.  On $y\ge\epsilon>0$ the convergence is uniform as all shifts
vanish; the integral over $0<y<\epsilon$ is uniformly small as
$\epsilon\downarrow0$.  Hence \eqref{gcx:after-subtraction} tends to
\[
 \int_0^\delta A_S(x;0)\lambda_S(\log x)\,\dd x-B_S(0).
\]
Here $P_S(x;0)$ is the Taylor polynomial of $q_S(x;0)$ through degree
$N_S-1$.  The two-endpoint finite part of each monomial
$x^{j-N_S}(\log x)^k$, $j<N_S$, is zero, so
\[
 -B_S(0)=\FP\int_0^\delta
          P_S(x;0)x^{-N_S}\lambda_S(\log x)\,\dd x.
\]
Thus the limit is precisely the directly merged finite part of the
subset contribution.  Summing the finitely many subsets and the outside
integral proves \eqref{gcx:limit}.  The proof did not fix an ordering
chamber; there are only finitely many choices for the minimum, and the
same dominating function applies to all of them.
\end{proof}

\begin{proof}[Proof of finite polylogarithmic reduction]
For distinct nodes, the rational principal part is
\begin{equation}\label{gcx:partial-fractions}
 \frac{P_S(x;\boldsymbol s)}{D_S(x;\boldsymbol s)}
  =\sum_{i\in S}\sum_{k=1}^{d_i}
       \frac{b_{S,i,k}(\boldsymbol s)}{(x+s_i)^k},
\end{equation}
where
\begin{equation}\label{gcx:b-coefficient}
 b_{S,i,k}=\frac1{(d_i-k)!}
 \left.\partial_x^{d_i-k}
 \frac{q_S(x;\boldsymbol s)}
      {\prod_{j\in S\setminus\{i\}}(x+s_j)^{d_j}}
 \right|_{x=-s_i}.
\end{equation}
These are finite expressions in the stated tail and weight data.
Translate to $y=x+b_S$.  All letters $-a_i$, where
$a_i=s_i-b_S\ge0$, lie on the nonpositive real axis.  For $a>0$,
\[
 \log(y+a)=\log a+G(-a;y),\qquad
 G(-a;y)=\int_0^y\frac{\dd t}{t+a},
\]
while $G(0;y)=\log y$ uses its usual tangential normalization.
The standard iterated-integral shuffle rule \cite{goncharov} expands each logarithmic product into words of
length at most $M$.  A simple pole in \eqref{gcx:partial-fractions}
prepends one letter and therefore gives weight at most $M+1$.

For a pole of order $k\ge2$, integration by parts gives a boundary term
of logarithmic degree at most $M$ and rational integrals of logarithmic
degree at most $M-1$.  Repeat after partial fractions; induction on $M$
terminates, even if a repeated pole is transferred to another letter.
This proves the same weight bound.  Lower and upper finite parts are the
constant terms of the resulting iterated integrals at $0$ and $\infty$;
the real integral in \eqref{gcx:T} fixes their branches without crossing
any interior letter.

When all $m_i=0$, every logarithmic polynomial is constant.  The upper
primitive has zero constant term, and evaluation at the lower endpoint
gives precisely the sums of logarithms and reciprocal powers in the
index-zero repeated-pole formula.  This proves exact compatibility with
the previous convention.
\end{proof}

\begin{remark}[What regularity has been proved]\label{gcx:regularity}
The result asserts continuity after the full subtraction at arbitrary
Stieltjes indices.  The residual logarithmic integral need not be
analytic in the separations.  When all $m_i=0$, its integrand is analytic
and the lower endpoint is an analytic coordinate on each ordering
chamber; then the earlier stronger analytic-extension statement is
recovered.  No differentiability of the higher-index remainder is used
to manufacture contact terms.
\end{remark}

% END sections/05_higher_collisions.tex


% BEGIN sections/06_collision_examples.tex
\subsection{One logarithmic index: an explicit dilogarithm kernel}
\label{gcx:one-log}

At total Stieltjes index one, the universal counterterms require only
ordinary dilogarithms.  Define, for $a,b\ge0$,
\[
 U(b,a)=\FP\int_0^\infty\frac{\log(y+a)}{y+b}\,\dd y.
\]
The lower finite part is used only when necessary, and the upper
subtraction is $\tfrac12\log^2R$.  The following exact table fixes every
branch by real positive arguments:
\begin{equation}\label{gcx:U-table}
\begin{array}{c|l}
 \text{condition}& U(b,a)\\ \hline
 a=b=0&0\\[2pt]
 b=0<a&\zeta(2)+\tfrac12\log^2a\\[2pt]
 a=0<b&-\zeta(2)-\tfrac12\log^2b\\[2pt]
 a=b>0&-\tfrac12\log^2a\\[2pt]
 a>b>0&\tfrac12\log^2(a-b)+\zeta(2)
       -\log(a-b)\log b+\Li_2\!\left(-\frac b{a-b}\right)\\[2pt]
 b>a>0&-\tfrac12\log^2(b-a)-\zeta(2)
       -\log a\log\!\left(\frac b{b-a}\right)
       -\Li_2\!\left(-\frac a{b-a}\right).
\end{array}
\end{equation}
The corresponding index-zero kernel is
\[
 V(b)=\FP\int_0^\infty\frac{\dd y}{y+b}
       =-\log b\quad(b>0),\qquad V(0)=0.
\]

\begin{proof}
If $a>b$, put $d=a-b$.  A primitive is
\[
 \log d\log(y+b)-\Li_2\!\left(-\frac{y+b}{d}\right).
\]
If $b>a$, put $d=b-a$.  A primitive is
\[
 \log(y+a)\log\!\left(1+\frac{y+a}{d}\right)
       +\Li_2\!\left(-\frac{y+a}{d}\right).
\]
Differentiation verifies both formulas on the real interval of
integration.  Evaluate their lower constant terms and use
\[
 \Li_2(-t)+\Li_2(-1/t)=-\zeta(2)-\tfrac12\log^2t,
 \qquad t>0,
\]
with the correct normalization $-\zeta(2)=-\pi^2/6$, to take the upper
constant.  The identity follows itself by differentiation and evaluation
at $t=1$, where $\Li_2(-1)=-\zeta(2)/2$.  Direct logarithmic primitives
handle $a=b$ and the endpoints $a=0$ or $b=0$.
\end{proof}

As a normalization check, the two-factor case with shifts $(0,a)$ gives
\[
 T_{(1,0)}(0,a)=\frac{\tfrac12\log^2a+\zeta(2)}a,\qquad
 T_{(0,1)}(0,a)=\frac{\log^2a+\zeta(2)}a.
\]
These are precisely the asymmetric binary collision polynomials already
proved in the manuscript.  Their unequal leading coefficients are
retained by the general construction.

\subsection{Three new mixed cubic collision constants}
\label{gcx:cubic}

For $0<\epsilon<1/4$ and $j=0,1,2$, set
\begin{equation}\label{gcx:Cj}
 C_j(\epsilon)=\FP\int_{\R/\Z}
  \gamma_1(\{x+j\epsilon\})
  \prod_{\substack{0\le k\le2\\ k\ne j}}
       \gamma_0(\{x+k\epsilon\})\,\dd x,
 \qquad
 Q=\FP\int_0^1\gamma_1(x)\gamma_0(x)^2\,\dd x.
\end{equation}
The same global constant $Q$ occurs in all three limits, because their
merged products agree pointwise before regularization.  Put
\[
 h(x)=\gamma_0(1+x),\qquad k(x)=\gamma_1(1+x),\qquad
 L=\log\epsilon,\quad \ell=\log2,\quad
 B(X)=\zeta(2)+\tfrac12X^2.
\]
Define
\begin{align}
 A_0(L)&=\tfrac14L^2-\tfrac12\ell L
                    +\tfrac12\zeta(2)-\tfrac14\ell^2,\notag\\
 A_1(L)&=\tfrac12L^2-\tfrac12\ell L+\tfrac14\zeta(2),
 \qquad A_2(L)=\tfrac12L^2.
 \label{gcx:Aj}
\end{align}

\begin{proposition}[Exact local counterterms for the three placements]
\label{gcx:three-T}
The counterterms of Definition~\ref{gcx:counterterm-definition} are
\begin{align}
 T_0(\epsilon)={}&\frac{A_0(L)}{\epsilon^2}
  +\frac{h(\epsilon)}\epsilon B(L)
  +\frac{h(-\epsilon)}{2\epsilon}B(L+\ell)
  +\frac{k(-2\epsilon)}\epsilon L,\label{gcx:T0}\\
 T_1(\epsilon)={}&\frac{A_1(L)}{\epsilon^2}
  +\frac{h(2\epsilon)}\epsilon B(L)
  +\frac{h(\epsilon)}{2\epsilon}L^2
  +\frac{k(-\epsilon)}{2\epsilon}(L+\ell)
  +\frac{h(-2\epsilon)}\epsilon B(L),\label{gcx:T1}\\
 T_2(\epsilon)={}&\frac{A_2(L)}{\epsilon^2}
  +\frac{k(\epsilon)}\epsilon L
  +\frac{h(\epsilon)}{2\epsilon}B(L+\ell)
  +\frac{h(-\epsilon)}{4\epsilon}(L+\ell)^2\notag\\
 &+\frac{h(-\epsilon)}\epsilon B(L)
  +\frac{h(-2\epsilon)}{2\epsilon}L^2.
 \label{gcx:T2}
\end{align}
In particular $C_j(\epsilon)-T_j(\epsilon)\to Q$ for each placement.
\end{proposition}

\begin{proof}
Only subsets of two and three poles contribute.  For the three-pole
subset the principal numerator is $1$, and
\[
 \frac1{x(x+\epsilon)(x+2\epsilon)}
 =\frac1{2\epsilon^2x}-\frac1{\epsilon^2(x+\epsilon)}
          +\frac1{2\epsilon^2(x+2\epsilon)}.
\]
Insert respectively $\log x$, $\log(x+\epsilon)$ and
$\log(x+2\epsilon)$, and apply \eqref{gcx:U-table}.  The answers are
$A_j(L)/\epsilon^2$.

For a pair of poles at shifts $a<b$, with omitted analytic tail $q$,
the principal part is
\[
 \frac{q(-a)}{(b-a)(x+a)}
    -\frac{q(-b)}{(b-a)(x+b)}.
\]
When the omitted index is one, use $V$; otherwise use $U$ with the
logarithm at the indicated pole.  The three pair subsets give exactly
the remaining terms of \eqref{gcx:T0}--\eqref{gcx:T2}.  Singleton
contributions vanish as already proved.  The limit is
Theorem~\ref{gcx:main}.
\end{proof}

For an expression having a power--log expansion through order zero,
$\CT_{\epsilon\downarrow0}$ denotes the coefficient of
$\epsilon^0(\log\epsilon)^0$.  It discards negative powers of $\epsilon$
and nonconstant powers of $\log\epsilon$, including the latter at power
$\epsilon^0$.

\begin{corollary}[The mixed finite anomalies]\label{gcx:CT-values}
With $\beta=\zeta(2)+\zeta'(2)$,
\begin{align}
 \CT_{\epsilon\downarrow0}C_0(\epsilon)
      &=Q-\tfrac12\zeta(2)^2+\tfrac14\zeta(2)\log^22,
 \label{gcx:CT0}\\
 \CT_{\epsilon\downarrow0}C_1(\epsilon)
      &=Q-\tfrac12\bigl(\zeta(2)+\zeta'(2)\bigr)\log2,
 \label{gcx:CT1}\\
 \CT_{\epsilon\downarrow0}C_2(\epsilon)
      &=Q+\tfrac12\zeta(2)^2.
 \label{gcx:CT2}
\end{align}
Consequently the global mixed constant cancels in the exact identity
\begin{equation}\label{gcx:CT-difference}
 \CT_{\epsilon\downarrow0}\bigl(C_2(\epsilon)-C_0(\epsilon)\bigr)
       =\zeta(2)^2-\tfrac14\zeta(2)\log^22.
\end{equation}
\end{corollary}

\begin{proof}
The necessary analytic expansions are
\[
 h(x)=\gamma-\zeta(2)x+O(x^2),\qquad
 k(x)=\gamma_1+\beta x+O(x^2).
\]
Here $\beta=k'(0)$ follows from \eqref{gcx:tail-derivative}; the derivative
of the Pochhammer factor contributes the term $\zeta(2)$.
Substitute into \eqref{gcx:T0}--\eqref{gcx:T2}.  Every omitted term is
$O(\epsilon(1+|\log\epsilon|)^2)=o(1)$, whereas
$C_j-T_j=Q+o(1)$.  Taking the constant coefficients proves all three
formulas.
\end{proof}

The first placement, with every divergent and finite logarithmic term
shown, is
\begin{align}
 C_0(\epsilon)={}&\frac1{\epsilon^2}
  \left\{\frac{L^2}{4}-\frac{\ell L}{2}
                  +\frac{\zeta(2)}2-\frac{\ell^2}{4}\right\}\notag\\
 &+\frac1\epsilon\left\{
       \frac{3\gamma L^2}{4}
       +\left(\frac{\gamma\ell}{2}+\gamma_1\right)L
       +\frac{3\gamma\zeta(2)}2+\frac{\gamma\ell^2}{4}\right\}\notag\\
 &-\frac{\zeta(2)L^2}{4}
       +\left(\frac{\zeta(2)\ell}{2}-2\beta\right)L
       +Q-\frac{\zeta(2)^2}{2}
                +\frac{\zeta(2)\ell^2}{4}+o(1).
 \label{gcx:C0-expansion}
\end{align}
The finite constants in \eqref{gcx:CT0}--\eqref{gcx:CT2} depend on which
factor carries the logarithmic Stieltjes singularity.  Retaining only
the leading inverse-power divergence would therefore miss the exact
answers.

\subsection{An ordinary integral for the remaining merged constant}

The constant $Q$ in \eqref{gcx:Cj} is not left without an independent
ordinary-integral definition.  Local Laurent subtraction gives
\begin{align}
 Q={}&-\frac14-2\gamma-\gamma_1\notag\\
 &+\int_0^1\Biggl\{
 \frac{2\bigl(h(x)-\gamma+\zeta(2)x\bigr)\log x}{x^2}
 +\frac{\bigl(h(x)^2-\gamma^2\bigr)\log x}{x}\notag\\
 &\hspace{17mm}+\frac{k(x)-\gamma_1-\beta x}{x^2}
 +\frac{2\bigl(h(x)k(x)-\gamma\gamma_1\bigr)}x
 +h(x)^2k(x)\Biggr\}\,\dd x.
 \label{gcx:Q-ordinary}
\end{align}
Each numerator has the stated vanishing order, so the integral is
absolutely convergent.  To verify the constant, expand
$(\log x/x+k(x))(1/x+h(x))^2$.  The pure singular terms have finite
integrals
\[
 \FP\int_0^1x^{-3}\log x\,\dd x=-\frac14,\quad
 \FP\int_0^1x^{-2}\log x\,\dd x=-1,\quad
 \FP\int_0^1x^{-2}\,\dd x=-1,
\]
while the $x^{-1}$ and $x^{-1}\log x$ terms have zero finite integral.
This proves \eqref{gcx:Q-ordinary}.  It supplies independent numerical
access to the global constant without converting the proof of the
collision identities into a numerical fit.

% END sections/06_collision_examples.tex


% BEGIN sections/07_audit_and_research.tex
\section{Audit, integration corrections, and further research}
\label{audit:section}

\subsection{What the audit changes}

No new false theorem in the focused established source material was
identified. The principal audit finding is a status issue: several
questions survive in earlier archives after later incoming reports have
answered them. Integration should reconcile the chronology before
assigning credit or declaring a question newly resolved.

\begin{longtable}{>{\raggedright\arraybackslash}p{0.29\textwidth}
                  >{\raggedright\arraybackslash}p{0.63\textwidth}}
\toprule
Earlier question or claim & Proposed integration treatment\\
\midrule
\endhead
Single-resolvent nonpositive powers &
The complete answer is already in \cite{repo:quartic}. The earlier
Barnes Question F5 is therefore resolved by that report. The new
theorem here concerns independently powered products and their
orientation-dependent endpoint poles.\\[5pt]
The fourth digamma power &
Retain the existing cubic and quartic bridges with credit to
\cite{repo:cubic,repo:barnes,repo:quartic}. Add the fifth theorem and the
all-power reduction as subsequent extensions.\\[5pt]
Pure pairs of even harmonic tails &
The question in the earlier parity and reflection reports is already
answered by \cite{repo:quartic}. It is not a new contribution of this
article. Products with further factors remain a separate problem.\\[5pt]
Higher-index collision remainder &
Replace the open-status entry for existence of a fully subtracted
total limit by Theorem~\ref{gcx:main}. State continuity explicitly.
Do not upgrade this to analyticity at positive Stieltjes indices or
to associativity of arbitrary partial-merger prescriptions.\\[5pt]
Canonical power polynomials &
Record Corollary~\ref{hp:cancellation} as elimination of the inherited
height-one coefficients $\eta,\theta$. The additional derivative
coordinate $\chi_4$ remains, so the stronger single-coordinate
all-power question is still open.\\[5pt]
Gaussian $S_6$ and revised $S_8$ &
Preserve their conjectural status. Keep the previously rejected
coefficient vector distinct from the revised candidate, and distinguish
period equality from membership in a specified formal relation system.\\
\bottomrule
\end{longtable}

The companion \src{integration/CLAIM_LEDGER.csv} records theorem labels,
source questions, and proposed status changes. The source snapshot
manifest gives the inspected revision and hashes of the current incoming
archives. The older reports are dependencies and provenance, not copied
into this deliverable as if they were new work.

\subsection{Exact corrections to tempting extrapolations}

Two natural extrapolations are disproved by the calculations here.
They should be excluded explicitly in any consolidated statement.

\paragraph{A value along a parameter path is not automatically the
coordinate finite part.}
For opposite resolvents, \eqref{mr:pathlimit} gives the three limits
$-2\pi\ii,-\pi\ii,0$ for the same specialized function $K-\pi\ii$.
Its coordinate finite part is $-\pi\ii$.
The proposed replacement is the full germ \eqref{mr:H}, retaining
$E=\sum\epsilon_j\lambda_j$ as well as $\Lambda=\sum\lambda_j$.
The phase numerator may be simplified using integrality of each exponent
only when that hypothesis actually holds. The theorem itself requires
only integrality of the total at the chosen resonance.

\paragraph{Taking the order-zero term of a collision does not commute
with merging its factors.}
The three constants \eqref{gcx:CT0}--\eqref{gcx:CT2} differ although
their directly merged integrand is the same
$\gamma_1\gamma_0^2$. The replacement is to subtract the entire
counterterm \eqref{gcx:T} before passing to the limit. Dividing the local
cutoff by a shrinking gap would alter this convention and its constants.
The terms involving
$\gamma_1'(1)=\zeta(2)+\zeta'(2)$ must retain both summands;
differentiation of the Pochhammer factor supplies the first one.

These counterexamples correct unsupported extensions. They are not
retroactive claims that the earlier single-resolvent or index-zero
theorems were false.

Two further proof requirements should accompany the integration.
Matching principal parts of $\psi^k$ leaves an entire-function
ambiguity; Lemma~\ref{hp:ML} supplies the contour argument that eliminates
it. Also, an exact finite harmonic identity is not the same statement as
its asymptotic regularization: use \eqref{hp:regularization} with finite
terminal harmonic sums first, and replace convergent terminal words by
their limits only after controlling their tails. These requirements are
met in the proofs above.

\subsection{Questions opened by the new results}

The following projects ask for exact identities or structural theorems.
They are questions rather than unproved assertions of reducibility.
Each has a concrete starting formula in this article.

\begin{question}[The derivative coordinate of weight four]
Can
\[
 \chi_4=2(\partial_1-\partial_2)Z_2(3,1)+\partial_1Z_2(2,2)
\]
be reduced to a smaller specified algebra of ordinary zeta derivatives
and height-one Laurent coefficients?
\end{question}
The first task is to differentiate identities valid with genuinely free
orders, rather than differentiate a duality known only at discrete
integer indices. The free-order stuffle identity used after
\eqref{hp:Er} is admissible in this sense. A useful result would be an
exact reduction or a formal separating functional for a precisely
enumerated family of differentiated relations. Such a formal obstruction
would still not prove arithmetic independence.

\begin{question}[The sixth power and canonical cancellation matrices]
Compute the complete sixth-power reduction and determine which
polynomials in $\psi$ eliminate its inherited lower-depth coordinates.
Does a uniform family of power polynomials remove all but the single
height-one coefficient at every degree?
\end{question}
The compiler of Theorem~\ref{hp:allpowers} already terminates at every
fixed degree and fixes its constants. What remains is to simplify the
resulting moving-order coordinates and compute the exact cancellation
matrix against lower powers. The coefficient of the top-depth term is
already known, namely $-k!$; it is the lower-depth relations that must be
understood. The surviving $\chi_4$ in the fifth case is a specific test.

\begin{question}[A common Hurwitz shift for all digamma powers]
Extend the reduction and the normalized primitive to
$\psi(x+a)^k$, with a free common shift $a$, and determine the finite
multiple-Hurwitz or Lerch family required.
\end{question}
At $a>0$ the unit-interval integral is ordinary; the principal poles of
the global meromorphic function have moved to $-a-\N_0$.
The contour subtraction and finite summation-by-parts method should be
carried out with the shifted lattice, keeping the endpoint normalization
explicit. Rational $a$ offers a precise next specialization to colored
multiple-zeta data.

\begin{question}[Higher mixed logarithmic resolvent products]
For arbitrary positive integers $r,s$, find finite multiple-polylogarithmic
expressions for the all-index transforms of $L_z^rL_w^s$.
What depth and moving-order derivative families are sufficient?
\end{question}
The two-factor formula \eqref{ml:Tu} and the ordinary cubic formula
\eqref{ml:triple} are initial cases. Start from the integer convolution
constraints of the Fourier coefficients, apply exact partial fractions,
and keep the complete local power--log polynomial in the Stieltjes
subtraction. A finite representation in a stated class is the target;
introducing an unrestricted new generating integral alone would not
provide that reduction.

\begin{question}[Classical multivariable forms of the Gamma spectrum]
For which exponent patterns and disk coordinates does
\eqref{mr:ck} collapse to a finite combination of Gauss, Appell,
Lauricella, or multiple-polylogarithmic functions?
\end{question}
The example \eqref{mr:gaussmean} proves that a Gauss function suffices
for one zero disk coordinate. The general Gamma convolution is already
explicit and normally convergent. A further result should produce a
smaller function class, a new functional equation, or an arithmetic
special-value identity, with all exceptional parameter restrictions
handled on the complete meromorphic expression.

\begin{question}[Real-axis boundary identities]
Determine the exact limits of the completed mixed resolvent transforms
as one or more nonreal parameters approach the real axis, including
coincident limiting parameters.
\end{question}
The sign $\epsilon_j$ then records which boundary value is taken, and
interior cotangent poles can appear. A correct statement must inventory
their local principal-value and contact contributions as well as the
endpoint terms already present in \eqref{mr:H}. This is an exact
boundary-value problem, not a request for numerical bounds.

\begin{question}[Regularity of the higher-index collision remainder]
Find the complete power--log expansion of $C-T$ near total collision,
and determine its differentiability and analytic subfamilies.
\end{question}
Theorem~\ref{gcx:main} proves continuity uniformly in the relative
gaps. It does not assert a convergent Taylor series. After the Hermite
subtraction the only obstacle is the explicit product of translated
logarithms. Resolving that remaining integral for one cluster and then
for nested clusters would give a controlled higher-order contact
calculus, including gap configurations not described by a single scale.

\begin{question}[The global mixed cubic constant]
Express
$Q=\FP\int_0^1\gamma_1(x)\gamma_0(x)^2\,\dd x$
in a specified multiple-Hurwitz or directional Tornheim derivative
family, and seek combinations that eliminate it.
\end{question}
Equation~\eqref{gcx:Q-ordinary} is an absolutely convergent representative.
The difference of collision constants already eliminates $Q$ exactly.
For a further arithmetic reduction, begin with a joint completed
three-factor Hurwitz generator and differentiate its orders before
specialization. A numerical relation for the displayed decimal would
only suggest a candidate.

\begin{question}[Compatibility of successive partial mergers]
Characterize partial-cluster subtraction rules whose iterated mergers
agree with the simultaneous full counterterm \eqref{gcx:T}.
Can a universal local adjustment make the operation associative?
\end{question}
The present result fixes the total limit even for hierarchical gaps.
It does not identify the finite constants produced by every possible
intermediate prescription. For three and four factors, compare two
merger orders symbolically using the explicit dilogarithm kernels,
then determine the local difference before proposing a general rule.

\begin{question}[Integer-frequency grids and their logarithmic scales]
Extend the geometric theorem to
$\prod_i\gamma_{m_i}^{(r_i)}(\{p_i x+s_i\})$ with positive integers
$p_i$, including analytic Fourier weights.
\end{question}
The analytic-weight portion at common unit frequency is proved here.
For a grid point with ambient displacement $t$, the singular factor
has amplitude
$p_i^{-r_i-1}t^{-r_i-1}L_{m_i,r_i}(\log t+\log p_i)$.
A complete theorem should inventory all coinciding grid points and
transport the cutoff in this ambient coordinate. The logarithms
$\log p_i$ are finite contacts that cannot be removed by an informal
rescaling.

\begin{question}[A common spectral and collision completion]
Construct a joint generator with both the complex resolvent-power
weight and moving periodic Stieltjes factors, and prove compatible
restriction formulas at intersecting spectral and geometric resonances.
\end{question}
The two ingredients are explicit here, but their direct combination
has new local products: a weighted collision point can approach an
endpoint where the weight is singular. The target is one full
subtraction rule that recovers \eqref{mr:H} and \eqref{gcx:T} in
their separate domains, with the order of all restrictions specified.

\begin{question}[A new exact relation for the frozen Gaussian targets]
Can an identity obtained from a harmonic, Stieltjes, or resolvent
transform supply a genuine functional relation for the frozen
$S_6$ or revised $S_8$ residual outside the previously obstructed
formal systems?
\end{question}
The source vectors and branch conventions must remain fixed. A useful
candidate should have a demonstrably new action on the existing
separating functional, or a proved analytic specialization beyond the
enumerated relation family. Neither additional digits nor an
uncompleted modular elimination establishes the period equality.
No such functional bridge has been proved in this article.

\begin{question}[Reusable formalization units]
Formalize the full-amplitude completion lemma, the all-power
Mittag--Leffler normalization, strict harmonic quasi-shuffles and their
asymptotic regularization, and the fixed-contour Hermite remainder.
\end{question}
These units separate the reusable analytic proof from finite symbolic
coefficient calculations. The supplied programs give regression
examples for the latter. Formal verification of the analytic lemmas
would support generated identities at many new indices; importing
Python outputs alone should not confer a proof-assistant status.

\subsection{Suggested order of work}

The most direct next exact calculations are the sixth-power normal
form, the first higher mixed logarithmic transform, and a joint
Hurwitz representation of the global mixed cubic constant. Their input
algebra and all local corrections are now explicit. The collision
regularity and merger-compatibility questions are the next structural
problems. The Gaussian conjectures require a new functional relation or
an exact period certificate beyond the relation systems already
analyzed in the source material.

% END sections/07_audit_and_research.tex


% BEGIN sections/08_verification.tex
\section{Verification, reproducibility, and the evidence boundary}
\label{verification:section}

\subsection{Exact calculations and independent numerical methods}

Every main assertion above is supported by an analytic proof or a finite
algebraic derivation. The code checks sensitive coefficients and supplies
independent numerical comparisons. Normal convergence, meromorphic
continuation, and the hierarchy-uniform collision limit are proved in
the article; they are not inferred from numerical agreement.

The harmonic compiler retains exact rational and symbolic coefficients.
Its sixty finite checks include the principal-part polynomial, the
printed fifth finite harmonic identity, the moving-order finite
lattices through degree six, and the separate regularization constants.
The classical convergent multiple-zeta identities used in the reduction
are stated explicitly in \eqref{hp:MZVtable}; they are mathematical
inputs, not theorems established by numerical testing.

The two fifth-power evaluators use different analytic representations.
One integrates the convergent local Taylor series of $(x\psi(x))^5$
on $(0,1/8)$ and uses ordinary quadrature on the remainder.
The other sums \eqref{hp:fifthseries}, using a finite polygamma
asymptotic tail represented by Hurwitz-zeta order derivatives.
At 70 working decimal digits their observed difference is
$3.3937772261\cdot10^{-49}$. Thus the recorded agreement is at the
48-decimal-place level; a 70-digit working precision is not itself
a 70-digit accuracy certificate.

The opposite-orientation tests verify the residue polynomial exactly
for $M=1,\ldots,9$, and the first logarithmic contact coefficients
symbolically. Direct beta quadrature at complex exponents has observed
residual $7.37\cdot10^{-83}$ at 80 working digits. Four transverse
rays and three curved paths recover their distinct uncompleted limits
and the common completed value.

The general Fourier convolution is also tested at three distinct
nonzero disk coordinates, with two upper orientations and one lower
orientation. The exponents are complex, one has negative real part,
and the total has real part $1.60$. At 80 working digits, the five modes
$k=-2,-1,0,1,2$, computed on the rectangle $0\le n,m\le160$, differ
from direct cotangent-resolvent quadratures by at most
$2.20\cdot10^{-80}$. A shorter rectangle through 120 gives observed
changes below $5.7\cdot10^{-63}$.
At exponents $(2,1)$ and zero disk coordinates the spectrum is finite.
The exponential moment is checked at four complex values, including
a reciprocal-Gamma zero, and the log-Gamma formula
\eqref{mr:loggammaexample} has observed residual below
$9.3\cdot10^{-84}$.

For the mixed-log transform, direct quadrature of the Bernoulli
polynomial $\zeta(-1,x)$ against the logarithmic product is compared
with separately evaluated multiple-polylogarithmic sums at three disk
coordinate pairs. One pair is complex and one is $(0,0)$.
The largest observed spectral residual is
$6.60\cdot10^{-82}$, and the mean comparisons are below
$6.4\cdot10^{-81}$ at 80 working digits.

There is also a direct singular-integral check of
\eqref{ml:specialStieltjes}. The left side uses coordinate subtraction
of its actual endpoint polynomial. Independently, $\eta$ is evaluated
by a subtracted positive-real Mellin integral. Both methods give
\[
 -6.1973961294626328099297651968714197451620285651100755907475\ldots.
\]
The observed residual rounds to zero at the 80-digit working precision;
the code checks a tolerance of $10^{-70}$. This is an observed
floating-point agreement, not an interval enclosure.

\subsection{Independent geometric-collision checks}

The script \src{verify_collision_kernels.py} compares all cases of the
dilogarithm table with independently subtracted ordinary integrals,
using 70 working decimal digits. Its eight comparisons have maximum
observed absolute residual $1.82\cdot10^{-71}$.
It also checks the three all-pole coefficients and the three finite
anomalies by exact symbolic algebra, treating
$\zeta(2),\zeta'(2),\gamma,\gamma_1,\log2$ as independent symbols.
Each symbolic residual is exactly zero.

The separate script \src{verify_higher_collision.py} evaluates the
circle integrals by subtracting their actual local singularities at
each of the three points. It does not use the proposed $T_j$ formulas
to define those integrals. It evaluates the directly merged ordinary
representative \eqref{gcx:Q-ordinary} independently, obtaining
\[
 Q=-2.87039237858604022876331895417289849797325401866\ldots.
\]
At five gaps from $0.03$ down to $0.0003$, the quantities $C_j-T_j-Q$
approach zero for every placement. Selected values are
\[
\begin{array}{c|rrr}
 \epsilon&j=0&j=1&j=2\\ \hline
 0.03&-0.0828460&0.184968&0.608974\\
 0.003&-0.00940356&0.0289980&0.0833617\\
 0.0003&-0.000951212&0.00376258&0.0107961
\end{array}
\]
These are finite-separation remainders, not equality errors expected
to vanish at each positive gap. The working precision is 65 decimal
digits. A convergent Hurwitz Taylor evaluator for $\gamma_1(1+x)$
was separately compared with the generalized Stieltjes function at
four points, with observed residual below $9\cdot10^{-67}$.

The numerical paths above illustrate the limit and check the explicit
formulas. Uniformity for arbitrarily hierarchical gaps follows from
the fixed-contour and dominated-convergence proof, not from sampling
a finite collection of geometric paths.

\subsection{Companion files and replay}

The ZIP includes the complete modular \LaTeX{} source, a flattened
standalone source, the compiled PDF, scripts and raw JSON records,
source hashes, and a claim-by-claim integration ledger.
The article can be built without the ProveIt checkout or network access.
The Python calculations require Python 3, SymPy, and mpmath; the recorded
environment uses Python 3.12.14, SymPy 1.14.0, and mpmath 1.3.0.

\begin{longtable}{>{\raggedright\arraybackslash}p{0.40\textwidth}
                  >{\raggedright\arraybackslash}p{0.52\textwidth}}
\toprule
File in \src{verification/} & Purpose\\
\midrule
\endhead
\src{harmonic_power_reduce.py} &
General exact finite harmonic compiler; accepts a power on the command line.\\[4pt]
\src{verify_harmonic_powers.py} &
Sixty exact finite identities and normal-form comparisons.\\[4pt]
\src{verify_fifth_numeric.py} &
Independent Taylor/quadrature and polygamma-series evaluations of $Q_5$.\\[4pt]
\src{check_mixed_orientation.py} &
Exact residue and logarithmic coefficients; beta and path diagnostics.\\[4pt]
\src{verify_multiple_resolvents.py} &
Three-parameter complex-power Gamma convolution against direct Fourier
quadrature, plus exponential and log-Gamma moment checks.\\[4pt]
\src{check_mixed_log.py} &
Ordinary mixed-log spectral transform and dilogarithmic mean.\\[4pt]
\src{verify_special_mixed_log.py} &
Independent Mellin and coordinate finite-part check of
\eqref{ml:specialStieltjes}.\\[4pt]
\src{verify_collision_kernels.py} &
Dilogarithm-kernel quadratures and exact cubic anomaly calculations.\\[4pt]
\src{verify_higher_collision.py} &
Direct singularity-subtracted circle quadratures and the merged constant.\\
\bottomrule
\end{longtable}

The top-level \src{run_checks.py} explains and runs the replay sequence;
the verification README distinguishes exact checks, observed numerical
residuals, and finite-separation remainders. The build scripts retain
the same article content in modular and standalone forms.


% END sections/08_verification.tex

\clearpage
\begin{thebibliography}{99}

% BEGIN references.tex
\bibitem{repo:main}
ProveIt Contributors, \emph{Polylogarithms and their Arithmetic Bridges},
four-volume manuscript and source status records, repository revision
\src{9035bb0eaaf70f3398711b489733093dcc52e946}.
\url{https://github.com/VladimirReshetnikov/ProveIt/tree/9035bb0eaaf70f3398711b489733093dcc52e946/Analysis/Polylogarithms/docs/manuscript}.
The companion incoming directory was inspected at the same revision.

\bibitem{repo:quartic}
ProveIt incoming report, \emph{Quartic Tails and Complex Resolvents},
11 October 2026.
Archive \src{ProveIt_Quartic_Tails_Resolvents_2026-10-11.zip},
at the revision in \cite{repo:main}.
Relevant members: \src{sections/quartic_bridge.tex},
\src{sections/complex_powers.tex}, and
\src{sections/research_questions.tex}, Questions 1--2 and 10--11.
\url{https://github.com/VladimirReshetnikov/ProveIt/blob/9035bb0eaaf70f3398711b489733093dcc52e946/docs/incoming/ProveIt_Quartic_Tails_Resolvents_2026-10-11.zip}.

\bibitem{repo:barnes}
ProveIt incoming report,
\emph{Barnes Resolvents and Higher Harmonic Identities},
11 October 2026.
Archive \src{ProveIt_Barnes_Resolvents_and_Harmonic_Identities_2026-10-11.zip},
at the revision in \cite{repo:main}.
Relevant members: \src{sections/05_quartic_harmonic.tex} and
\src{sections/07_audit_research.tex}, Questions F3 and F5.
\url{https://github.com/VladimirReshetnikov/ProveIt/blob/9035bb0eaaf70f3398711b489733093dcc52e946/docs/incoming/ProveIt_Barnes_Resolvents_and_Harmonic_Identities_2026-10-11.zip}.

\bibitem{repo:cubic}
ProveIt incoming report,
\emph{Cubic Tornheim Derivatives and Harmonic Stieltjes Constants},
11 October 2026.
Archive \src{ProveIt_Cubic_Harmonic_Mixed_Orders_2026-10-11.zip},
at the revision in \cite{repo:main}.
Relevant member: \src{sections/02_cubic_harmonic.tex}.
\url{https://github.com/VladimirReshetnikov/ProveIt/blob/9035bb0eaaf70f3398711b489733093dcc52e946/docs/incoming/ProveIt_Cubic_Harmonic_Mixed_Orders_2026-10-11.zip}.

\bibitem{repo:reflection}
ProveIt incoming report,
\emph{Higher Reflection Identities and Resonant Zeta Calculus},
11 October 2026.
Archive \src{ProveIt_Higher_Reflection_Identities_2026-10-11.zip},
at the revision in \cite{repo:main}.
Relevant members: \src{sections/04_collision.tex}, Theorems
\src{gc:main} and \src{gc:derivative}, and
\src{sections/06_audit_research.tex}, the higher Stieltjes geometric
collision and integer-frequency/weight questions.
\url{https://github.com/VladimirReshetnikov/ProveIt/blob/9035bb0eaaf70f3398711b489733093dcc52e946/docs/incoming/ProveIt_Higher_Reflection_Identities_2026-10-11.zip}.

\bibitem{repo:exact}
ProveIt incoming report,
\emph{Exact Resonant Identities Beyond the Diagonal},
11 October 2026.
Archive \src{ProveIt_Exact_Resonant_Identities_2026-10-11.zip},
at the revision in \cite{repo:main}.
Relevant members: source audit, research-status section, and the
specified Gaussian/Cayley relation certificates.
\url{https://github.com/VladimirReshetnikov/ProveIt/blob/9035bb0eaaf70f3398711b489733093dcc52e946/docs/incoming/ProveIt_Exact_Resonant_Identities_2026-10-11.zip}.

\bibitem{dlmf:beta}
NIST Digital Library of Mathematical Functions,
\S5.12, \emph{Beta Function}, especially Eq.~5.12.6,
the Fourier beta integral.
\url{https://dlmf.nist.gov/5.12.E6}.

\bibitem{dlmf:hurwitz}
NIST Digital Library of Mathematical Functions,
\S25.11, \emph{Hurwitz Zeta Function},
including the recurrence, Fourier series, argument derivatives,
and the derivative at order zero.
\url{https://dlmf.nist.gov/25.11}.

\bibitem{dlmf:gamma}
NIST Digital Library of Mathematical Functions,
\S5.7, \emph{Series Expansions}, Gamma and digamma series.
\url{https://dlmf.nist.gov/5.7}.

\bibitem{dlmf:polygamma}
NIST Digital Library of Mathematical Functions,
\S5.15, \emph{Polygamma Functions}, especially the positive-integer
values and recurrence.
\url{https://dlmf.nist.gov/5.15}.

\bibitem{dlmf:psi}
NIST Digital Library of Mathematical Functions,
\S\S5.5, 5.7, 5.11, 25.2, and 25.16,
digamma recurrence and reflection, local and asymptotic expansions,
the Stieltjes convention, and Euler sums.
\url{https://dlmf.nist.gov/5.5},
\url{https://dlmf.nist.gov/5.11},
\url{https://dlmf.nist.gov/25.16}.

\bibitem{hoffman}
M.~E.~Hoffman, \emph{Algebraic aspects of multiple zeta values},
in \emph{Zeta Functions, Topology and Quantum Physics},
T.~Aoki et al. (eds.), Springer, 2005, pp.~51--74.
Preprint arXiv:math/0309425v2, 2004.
\url{https://arxiv.org/abs/math/0309425}.

\bibitem{mezo}
I.~Mez\H{o}, \emph{Nonlinear Euler sums},
Pacific Journal of Mathematics \textbf{272} (2014), no.~1, 201--226.
DOI: \src{10.2140/pjm.2014.272.201}.
\url{https://msp.org/pjm/2014/272-1/pjm-v272-n1-p09-s.pdf}.

\bibitem{goncharov}
A.~B.~Goncharov, \emph{Multiple polylogarithms and mixed Tate motives},
arXiv:math/0103059v4, 2001.
\url{https://arxiv.org/abs/math/0103059}.

% END references.tex

\end{thebibliography}
\end{document}
